\documentclass[12pt,a4paper]{report}

% -------------------------
% Encoding and fonts
% -------------------------
\usepackage[utf8]{inputenc}
\usepackage[T2A]{fontenc}
\usepackage[russian,english]{babel}   % English main, Russian secondary

% -------------------------
% Page layout and spacing
% -------------------------
\usepackage{geometry}
\geometry{margin=1in}
\usepackage{setspace}
\doublespacing

% -------------------------
% Mathematics, tables, and figures
% -------------------------
\usepackage{graphicx}
\usepackage{pdfpages}
\usepackage{xcolor}
\usepackage{caption}
\usepackage{amsmath}
\usepackage{amssymb}
\usepackage{booktabs}
\usepackage{tabularx}
\usepackage{float}
\usepackage{placeins}
\usepackage{array}
\usepackage{pdflscape}
\usepackage{makecell}
\usepackage{longtable}
\newcolumntype{L}[1]{>{\raggedright\arraybackslash}p{#1}}
\newcolumntype{C}[1]{>{\centering\arraybackslash}p{#1}}

% -------------------------
% TikZ, if needed for diagrams
% -------------------------
\usepackage{tikz}
\usetikzlibrary{arrows.meta,positioning,fit,calc}

% -------------------------
% References and hyperlinks
% -------------------------
\usepackage{natbib}
\bibliographystyle{plainnat}
\usepackage[hidelinks]{hyperref}
% -------------------------
% Safe inclusion helpers for notebook-exported figures/tables
% If a file is missing, the thesis still compiles and shows a visible placeholder.
% -------------------------
\newcommand{\safefigure}[2][]{%
  \IfFileExists{#2}{\includegraphics[#1]{#2}}{%
  \fbox{\parbox{0.85\textwidth}{\centering Missing file: \texttt{\detokenize{#2}}}}}%
}

\newcommand{\safetablepdf}[2][]{%
  \IfFileExists{#2}{\includegraphics[#1]{#2}}{%
  \fbox{\parbox{0.85\textwidth}{\centering Missing table file: \texttt{\detokenize{#2}}}}}%
}

% -------------------------
% Custom title page (HSE Faculty of Computer Science)
% -------------------------
\makeatletter
\renewcommand{\maketitle}{%
\begin{titlepage}
\centering

% Top: university info
{\scshape\large
FEDERAL STATE AUTONOMOUS EDUCATIONAL INSTITUTION \\
FOR HIGHER EDUCATION \\
NATIONAL RESEARCH UNIVERSITY \\
``HIGHER SCHOOL OF ECONOMICS''
\par}

\vspace{1cm}

% Faculty
{\large Faculty of Computer Science\par}

\vspace{1cm}

% Author name
{\Large \textbf{Muhammad Suliman}\par}

\vspace{1cm}

% Russian title (optional, can be left blank or same as English)
{\large Russian title: \textit{``Экономические и экологические потери от лесных пожаров с использованием машинного обучения''}\par}

\vspace{0.5cm}

% English title
{\Large \textbf{Economic and Ecological Losses from Wildfires using Machine Learning}\par}

\vspace{1cm}

% Qualification line
{\large Qualification paper -- Master of Science Dissertation\par}

\vspace{0.5cm}

% Field of study
{\large Field of study 01.04.02 ``Applied Mathematics and Informatics''\par}
{\large Program: ``Data Science''\par}

\vspace{0.5cm}

\begin{flushright}
\begin{tabularx}{0.5\textwidth}{X}
Student: \textbf{Muhammad Suliman} \\[0.5cm]
Supervisor: \textbf{Ramon Antonio Rodriges Zalipynis} \\
\end{tabularx}
\end{flushright}

\vfill

% Bottom: city and year
{\large Moscow, 2026\par}

\end{titlepage}
}
\makeatother

\begin{document}

% -------------------------
% Front matter
% -------------------------
\pagenumbering{roman}

\maketitle
\chapter*{Abstract}
\addcontentsline{toc}{chapter}{Abstract}

Wildfires are increasingly important socioecological hazards whose consequences extend beyond fire occurrence and burned area. Although recent machine-learning research has improved wildfire prediction, fewer studies connect wildfire hazard information with economic and ecological consequence assessment. This thesis develops a county-level wildfire risk prediction and assessment pipeline using publicly accessible datasets. The main empirical target is wildfire Expected Annual Loss from the FEMA National Risk Index, which provides a monetary county-level loss indicator. Baseline exposure, value, vulnerability, and resilience predictors are combined with environmental variables derived from topography, land cover, and warm-season climate summaries.

Several regression models are evaluated, including Linear Regression, Ridge Regression, Random Forest, and Gradient Boosting. The results show that the county-only baseline provides meaningful predictive signal, while environmental augmentation substantially improves performance. The strongest results are obtained with the combined Random Forest model, and state-wise grouped validation confirms that the model retains useful predictive power under stricter geographic testing. The pipeline also produces county-level risk-assessment outputs by converting predicted log-transformed wildfire EAL back into monetary units and ranking counties by predicted wildfire loss.

Because the thesis title also addresses ecological consequences, an exploratory ecological risk-screening layer is added. This layer does not estimate ecological losses directly; instead, it combines predicted wildfire-risk percentile with natural vegetation exposure to identify counties where high monetary wildfire risk overlaps with high forest, shrub, and grass coverage.

The novel contribution of the thesis is the general consequence-oriented modeling framework rather than the specific territory of application. Unlike much previous wildfire-machine-learning work that focuses mainly on fire occurrence, susceptibility, spread, or burned area, this thesis connects wildfire-relevant environmental predictors with exposure, vulnerability, resilience, and a monetary loss-oriented target. The new element is the integration of county-level wildfire Expected Annual Loss modeling, environmental augmentation, state-wise geographic validation, and an exploratory ecological risk-screening layer within one reproducible workflow. The main insights gained are that environmental predictors substantially improve wildfire-loss prediction, that geographic validation gives a more realistic estimate of transferability, and that monetary wildfire risk can be complemented by ecological exposure screening.

In addition, supporting results from the An Giang, Vietnam wildfire prediction study are reported. Part of this work was successfully presented at the 4th International Conference on Distributed Computing and High Performance Computing (DCHPC 2026), and the associated paper is expected to appear as a Web of Science/Scopus paper. The conference certificate is included in Appendix~\ref{app:conference_certificate}.

\chapter*{Description of the Use of a Generative Model}
\addcontentsline{toc}{chapter}{Description of the Use of a Generative Model}

During the preparation of this thesis, generative artificial intelligence tools were used only as auxiliary support for improving academic style, grammar, clarity, and wording. The research design, data preparation, experiments, numerical results, interpretation, conclusions, and final decisions were completed by the author in tight collaboration with the supervisor.

\tableofcontents
\listoffigures
\listoftables

\clearpage

% -------------------------
% Main matter
% -------------------------
\pagenumbering{arabic}

\chapter{Introduction}
\label{chap:introduction}

Wildfires are becoming an increasingly significant source of economic, ecological, and public-health disruption in many regions of the world. In recent years, machine learning and deep learning have substantially improved the ability to predict wildfire occurrence, susceptibility, spread, and short-term fire danger. However, a major part of the literature still treats wildfire hazard prediction as an endpoint rather than as an input to broader consequence analysis. As a result, there remains a gap between wildfire prediction and wildfire-loss assessment, especially when the objective is to estimate impacts in monetary units or to connect fire-related conditions with broader ecological consequences.

This thesis is motivated by that gap. The central idea is that wildfire risk should be viewed not only as the probability of fire occurrence, but also as a function of what is exposed, how vulnerable it is, and what losses may follow. In this sense, modern predictive modeling should support not only hazard estimation but also consequence-oriented analysis. This perspective is increasingly emphasized in recent literature on wildfire expected annual loss, economy-wide disruption, smoke-related health burdens, and ecological vulnerability \citep{asif2026supranational,wang2021economicfootprint,law2025anthropogenic,arrogantefunes2024ecological}.

The present work focuses on county-level wildfire-loss modeling using publicly accessible datasets. As a first experimental direction, wildfire expected annual loss from the FEMA National Risk Index is used as a monetary target variable. This target is then modeled using county-level demographic, exposure, and resilience indicators, and later augmented with environmental predictors derived from publicly accessible geospatial datasets. In this way, the thesis develops a scalable baseline that connects wildfire-relevant environmental conditions with a consequence-oriented loss metric.

The main objectives of the thesis are:
\begin{enumerate}
    \item to review recent literature at the intersection of wildfire prediction, wildfire-loss assessment, and emerging AI methods;
    \item to construct a county-level wildfire-loss dataset in monetary units using accessible public data;
    \item to evaluate whether county-level wildfire expected annual loss can be predicted from exposure, vulnerability, and environmental variables;
    \item to assess how much environmental augmentation improves the predictive performance and geographic robustness of the resulting models.
\end{enumerate}

The novelty of this thesis is not limited to a particular study territory. Instead, the thesis proposes a general consequence-oriented wildfire-risk modeling workflow that can be adapted to other regions when comparable hazard, exposure, environmental, vulnerability, and loss-related data are available. The empirical implementation uses FEMA NRI wildfire Expected Annual Loss because it provides a standardized public monetary loss target, but the methodological logic is broader than this specific dataset.

Unlike much previous wildfire-machine-learning work that treats wildfire occurrence, susceptibility, fire danger, or burned area as the final output, this thesis uses wildfire prediction as a step toward loss-oriented risk assessment. The new contribution is to connect environmental wildfire drivers, exposure and value indicators, social vulnerability, community resilience, and monetary expected loss in one machine-learning pipeline.

The main contributions of the thesis are as follows:

\begin{enumerate}
    \item It provides a novel consequence-oriented framing of wildfire-risk modeling by shifting the focus from hazard prediction alone toward wildfire Expected Annual Loss estimation.
    
    \item It develops a new end-to-end machine-learning workflow that combines public risk indicators, exposure variables, vulnerability and resilience measures, topography, land cover, and climate predictors.
    
    \item Unlike models based only on socioeconomic or exposure variables, the thesis demonstrates that environmental augmentation substantially improves wildfire-loss prediction.
    
    \item It provides new insight into geographic transferability by comparing conventional random validation with state-wise grouped validation.
    
    \item It adds a novel exploratory ecological risk-screening layer that links predicted wildfire risk with natural vegetation exposure, without claiming to directly estimate ecological loss.
    
    \item It shows that the same general workflow can be extended beyond the current empirical territory to other regions, scales, and loss definitions, including event-level losses, asset-level exposure, ecological vulnerability, and post-fire recovery indicators.
\end{enumerate}

The experimental contribution of the thesis is a county-level wildfire Expected Annual Loss modeling pipeline built from FEMA National Risk Index data, followed by environmental augmentation using topography, land cover, and warm-season climate predictors. The results show that county-level wildfire loss can be modeled with meaningful predictive power and that environmental variables substantially improve both predictive performance and geographic transferability.

The implemented empirical workflow is intentionally county-level and national-scale, rather than event-level, pixel-level, or parcel-level. This choice enables reproducible modeling with accessible public data, while fine-scale wildfire tensor construction is retained as a future extension in Appendix~\ref{app:creek_fire_pipeline}.

Although the broader thesis theme includes both economic and ecological wildfire consequences, the implemented empirical pipeline focuses primarily on monetary wildfire Expected Annual Loss because FEMA NRI provides a standardized, publicly accessible county-level loss target. Ecological losses are not directly modeled as a separate damage target; however, the thesis adds an exploratory ecological risk-screening layer using natural vegetation exposure and predicted wildfire risk. This proxy-based analysis should be interpreted as ecological exposure/risk screening rather than direct ecological-loss estimation.

In addition to the main FEMA/NRI county-level pipeline, the thesis also briefly reports key results from a previous Vietnam case study on wildfire risk prediction in An Giang Province. That study used a compact deep neural network with multi-source environmental predictors, fire-growth feature engineering, threshold optimization, and SHAP-based interpretability. It is included as supporting previous work because it demonstrates the author's earlier experience with fine-scale wildfire prediction, while the main empirical contribution of this thesis remains the FEMA/NRI county-level wildfire Expected Annual Loss assessment pipeline.

Part of this work was successfully presented at the 4th International Conference on Distributed Computing and High Performance Computing (DCHPC 2026), held on 10--11 May 2026 in Tehran, Iran. The presented paper was entitled ``Mapping Forest Fire Risks Using Deep Learning: A Case Study in An Giang Province, Vietnam.'' The corresponding conference certificate is included in Appendix~\ref{app:conference_certificate}. The associated paper is expected to appear as a Web of Science/Scopus paper.

The remainder of the thesis is organized as follows. Chapter~\ref{chap:litreview} reviews the recent literature on wildfire risk, loss, ecological vulnerability, and emerging AI methods. Chapter~\ref{chap:methodology} describes the dataset construction process, predictor engineering, modeling workflow, and evaluation strategy. Chapter~\ref{chap:results} presents the experimental results and risk assessment. Chapter~\ref{chap:discussion} discusses the implications and limitations of the findings. Finally, Chapter~\ref{chap:conclusion} concludes the thesis and outlines directions for future work.
\chapter{Literature Review}
\label{chap:litreview}

\section{Introduction}

Wildfires are increasingly recognized as complex socioecological hazards whose consequences extend far beyond fire occurrence itself. Recent literature shows that machine learning, deep learning, remote sensing, geospatial data fusion, and multimodal artificial intelligence have substantially improved the ability to predict wildfire ignition, susceptibility, spread, fire danger, active-fire detection, and short-term fire behavior. At the same time, a smaller but rapidly growing body of work connects wildfire prediction to broader consequence assessment, including suppression expenditures, direct asset damage, expected annual loss, supply-chain disruption, smoke-related health burdens, and ecological degradation. This chapter reviews influential and recent literature across these connected areas and positions the present thesis at the intersection of wildfire-risk prediction and wildfire-loss assessment.

The review is organized around several connected themes. First, it clarifies the conceptual foundations of wildfire risk, hazard, exposure, vulnerability, resilience, and loss. Second, it discusses foundational fire-behavior, fire-regime, and wildfire-climate studies that motivate the use of environmental predictors. Third, it reviews data foundations, including remote sensing, meteorological datasets, land-cover products, and wildfire benchmark datasets. Fourth, it examines machine-learning and deep-learning approaches for wildfire prediction and monitoring. Fifth, it reviews the literature that moves from wildfire prediction to economic loss assessment, including expected annual loss, suppression costs, property damage, supply-chain disruption, economy-wide impacts, and smoke-related mortality. Sixth, it discusses agricultural and natural-hazard loss modeling as a methodological bridge. Seventh, it reviews spatial transferability and geographic validation, ecological losses, and emerging LLM/foundation-model approaches. Finally, it identifies the research gap addressed by this thesis.

\section{Wildfire Risk, Hazard, Vulnerability, and Loss}

A recurring issue in the wildfire literature is the inconsistent use of the terms \emph{risk}, \emph{hazard}, \emph{danger}, \emph{exposure}, \emph{vulnerability}, \emph{resilience}, and \emph{loss}. Recent wildfire-risk reviews emphasize that wildfire risk cannot be reduced to fire occurrence alone, because the impact of a wildfire depends not only on where and when fire occurs, but also on what assets, populations, infrastructure, and ecosystems are exposed and how susceptible they are to damage \citep{xu2024survey,xu2025deeplearning,jodhani2026ai,caron2025ai,ejaz2025mlpipeline}. In this sense, wildfire hazard generally refers to the likelihood, intensity, or physical potential of fire occurrence, while vulnerability captures the susceptibility of exposed systems to damage and their capacity to recover \citep{arrogantefunes2024ecological,asif2026supranational}.

This distinction is central to the present thesis because much of the machine-learning wildfire literature remains hazard-centric. Many studies stop at susceptibility maps, ignition probability, spread forecasts, or fire-danger scores, while fewer translate predictive outputs into explicit estimates of economic or ecological loss \citep{asif2026supranational,caron2025ai}. Wildfire-loss assessment therefore requires a broader framework in which hazard, exposure, vulnerability, and resilience are explicitly combined to estimate expected consequences \citep{asif2026supranational,ma2022supplychain,monge2023probabilistic,thompson2011risk,calkin2014risk}.

Influential fire-regime literature also shows that wildfire risk is shaped by the interaction of climate, vegetation, land use, and human systems. Bowman et al.\ describe fire as a central Earth-system process, while later work emphasizes the human dimension of global fire regimes, showing that fire is not only a biophysical process but also a product of land management, settlement patterns, ignition sources, and suppression practices \citep{bowman2009fire,bowman2011human}. Moritz et al.\ argue that societies need to learn to coexist with wildfire by shifting from a purely suppression-oriented paradigm to a risk-reduction and resilience perspective \citep{moritz2014learning}. These studies support the view that wildfire loss cannot be explained by fire occurrence alone; it depends on the coupled human--environment system in which fire occurs.

A similar conceptual separation is needed when distinguishing economic and ecological losses. Economic losses include suppression expenditures, damage to property and infrastructure, indirect business interruptions, supply-chain disruptions, public-health costs linked to smoke exposure, and broader macroeconomic impacts \citep{huang2024suppression,wang2021economicfootprint,monge2023probabilistic,law2025anthropogenic,qiu2025wildfiresmoke}. Ecological losses include biodiversity reduction, vegetation degradation, loss of ecosystem functions, reduced ecosystem-service capacity, reduced resilience, and long recovery periods after severe fires \citep{arrogantefunes2024ecological,rocesdiaz2022ecosystemservices,pacheco2023ecosystemservices,taboada2021ecosystemdelivery}. A complete wildfire-loss framework should therefore not rely on burned area alone as a proxy for impact, but should account for the value of exposed systems, the severity of damage, and recovery dynamics.

\section{Foundational Fire-Behavior, Climate, and Fire-Regime Literature}

Wildfire modeling has long been supported by physical and empirical fire-behavior models. Rothermel's surface fire-spread model remains one of the foundational contributions for predicting fire spread in wildland fuels \citep{rothermel1972spread}. The Canadian Forest Fire Weather Index system provides another influential operational framework for translating weather variables into fire-danger indicators \citep{vanwagner1987fwi}. FARSITE extended fire-behavior modeling into spatial simulation, while fuel-model references by Anderson and Scott and Burgan formalized important fuel categories used in fire-behavior estimation \citep{finney1998farsite,anderson1982fuelmodels,scott2005fuelmodels}. BehavePlus and related fire-behavior systems further demonstrate the long-standing operational importance of structured fuel, weather, and topographic inputs \citep{andrews2014behaveplus}.

Although this thesis uses machine learning rather than a physics-based fire-spread simulator, these foundational works are still important because they clarify why terrain, fuels, wind, moisture, and weather are central wildfire predictors. They also provide conceptual continuity between classical fire-science variables and modern data-driven feature engineering. Earlier probability-based and GIS/remote-sensing wildfire-risk frameworks similarly demonstrate that fire risk has long been modeled as a combination of hazard, environmental conditions, exposure, and values at risk \citep{preisler2004forecasting,chuvieco2010framework,ager2010wildfire,thompson2011risk}.

A large body of influential wildfire-climate research links changing fire regimes to warming temperatures, increasing atmospheric dryness, and longer fire-weather seasons. Westerling et al.\ showed that warming and earlier spring conditions contributed to increased western U.S. forest wildfire activity, establishing one of the key empirical links between climate change and wildfire activity \citep{westerling2006warming}. Jolly et al.\ demonstrated that global fire-weather seasons lengthened across many regions between 1979 and 2013, indicating that climate change can expand the temporal window in which large fires are possible \citep{jolly2015climate}. Flannigan et al.\ and Moritz et al.\ further discuss how changing climate can alter global fire activity, while Barbero et al.\ show increasing potential for very large fires in the contiguous United States \citep{flannigan2009climate,moritz2012climate,barbero2015climate}.

Environmental controls on wildfire operate across multiple spatial and temporal scales. Parisien and Moritz analyze environmental controls on wildfire distribution at multiple scales, while Littell et al.\ connect climate variability to area burned across western U.S. ecoprovinces \citep{parisien2009environmental,littell2009climate}. Krawchuk et al.\ introduce the broader concept of global pyrogeography, and Abatzoglou et al.\ show that climate--fire relationships vary across regions \citep{krawchuk2009global,abatzoglou2018global}. Human influence is also central: Syphard et al.\ examine human influence on California fire regimes, and Balch et al.\ show that human-started wildfires expand the fire niche across the United States \citep{syphard2007human,balch2017human}. In tropical forest systems, Arag{\~a}o and Shimabukuro show that fire incidence in Amazonian forests is closely connected to land-use change and has important implications for forest conservation and REDD policy \citep{aragao2010incidence}. Dennison et al.\ document large-wildfire trends in the western United States, reinforcing the importance of long-term fire-regime changes \citep{dennison2014large}.

Recent work highlights atmospheric aridity and moisture availability. Abatzoglou and Williams show that anthropogenic climate change increased fuel aridity across western U.S. forests and contributed substantially to increased forest fire area \citep{abatzoglou2016anthropogenic}. Williams et al.\ further show that increased atmospheric aridity played a major role in observed increases in California wildfire activity \citep{williams2019observed}. Trugman et al.\ emphasize the role of moisture availability and timing, while Holden et al.\ show that soil moisture and plant hydraulic stress can be strong predictors of fire-spread potential \citep{trugman2018flash,holden2025soilmoisture}. These studies provide important physical justification for including climate variables such as precipitation, maximum temperature, wind speed, and vapor pressure deficit in wildfire-risk and wildfire-loss models.

Global fire and emissions studies also provide context for wildfire impacts. Global fire-emissions datasets and fire radiative power methods show how remote sensing can quantify combustion, emissions, and biomass burning activity at large scales \citep{randerson2012gfed,van_der_werf2017global,wooster2005fire}. The State of Wildfires 2024--2025 report provides a recent global assessment of extreme wildfire episodes, exposure, emissions, and climate drivers \citep{kelley2025state}. Together, foundational fire-behavior, climate, and fire-regime studies provide the scientific basis for the environmental component of the present thesis.

\section{Data Foundations for Wildfire Modeling}

The recent wildfire literature is strongly data-driven. Review papers consistently show that wildfire-risk models are built from a combination of remote-sensing imagery, meteorological and climate data, vegetation and fuels information, topography, hydrology, historical fire records, and socioeconomic variables \citep{jain2020review,xu2025deeplearning,xu2024survey,ejaz2025mlpipeline,jodhani2026ai,liu2025advancements}. Satellite platforms such as Landsat, Sentinel, MODIS, VIIRS, GOES, LiDAR, SAR, and UAV-based systems have become central to wildfire research because they enable broad spatial coverage, repeated observations, and the derivation of key indices such as NDVI, NDWI, NBR, land surface temperature, canopy structure, burn scars, and fuel moisture \citep{jodhani2026ai,xu2025deeplearning}.

Remote-sensing fire products are especially important for active fire detection, burned-area mapping, and burn-severity assessment. The Monitoring Trends in Burn Severity (MTBS) project provides an influential reference for tracking burn severity over time \citep{eidenshink2007mtbs}. The MODIS active fire product and the VIIRS 375 m active fire product are widely used for active fire detection and near-real-time fire monitoring \citep{giglio2016modis,schroeder2014viirs}. Burn severity and spectral indices are also important for post-fire impact analysis: Key and Benson formalized the Composite Burn Index and normalized burn ratio approach, while Miller and Thode developed the relative differenced normalized burn ratio for heterogeneous landscapes \citep{key2006firemon,miller2007severity}.

Recent benchmark-oriented work highlights that data-source selection is itself a major methodological decision. Xu et al.\ introduce Sen2Fire, a Sentinel-based wildfire detection benchmark that combines Sentinel-2 multispectral imagery with Sentinel-5P aerosol information and evaluates the role of spectral bands and indices such as NBR and NDVI \citep{xu2024sen2fire}. Karlsson et al.\ compare MODIS and VIIRS products for next-day wildfire predictability and show that the choice of satellite fire-mask product can substantially affect learning behavior and prediction quality \citep{karlsson2025modisviirs}. Huot et al.\ provide a machine-learning dataset for next-day wildfire spread prediction from remote-sensing data, illustrating how wildfire modeling is moving toward benchmarked spatiotemporal learning tasks \citep{huot2022nextday}.

Environmental datasets also play a major role. The gridMET dataset provides gridded daily surface meteorological variables designed for ecological and environmental applications \citep{abatzoglou2013gridmet}. The National Land Cover Database provides Landsat-based land-cover products at 30 m resolution, making it suitable for deriving county-level land-cover fractions such as forest, shrub, grass, and developed land \citep{dewitz2021nlcd}. Vegetation and fuel-mapping work further shows why fuel and land-cover information are essential for fire-risk assessment and fire-management planning \citep{keane2001mapping}. These datasets are directly relevant to the present thesis because they support scalable construction of county-level environmental predictors for wildfire expected annual loss modeling.

Despite these advances, the literature identifies several data limitations. Cross-ecosystem transferability remains weak because models trained in one biome can perform poorly in another. Real-time integration remains limited in many studies, which often rely on static or historical datasets rather than dynamically updated data streams. Uncertainty quantification is also underdeveloped, especially for operational systems where probabilistic outputs would often be more useful than single deterministic scores \citep{jodhani2026ai,caron2025ai}. Wildfire datasets also frequently exhibit class imbalance, spatial autocorrelation, and inconsistent labeling across data products \citep{xu2024survey,caron2026multitarget,karlsson2025modisviirs}.

A second challenge is that data sufficient for wildfire prediction are not always sufficient for wildfire-loss estimation. Loss-focused studies require additional exposure and vulnerability layers, such as buildings, transport links, forests, energy infrastructure, demographic profiles, agricultural values, and other economic-value indicators \citep{asif2026supranational,ma2022supplychain,wang2021economicfootprint}. The transition from wildfire prediction to wildfire-loss modeling is therefore both a methodological shift and a data-expansion problem.

\section{Machine Learning and Deep Learning for Wildfire Prediction}

\subsection{Traditional Machine Learning}

Traditional machine-learning models, including logistic regression, random forests, support vector machines, gradient boosting, and related ensemble methods, are widely used for wildfire ignition, susceptibility, and danger mapping \citep{jain2020review,sayad2019predictive,liu2025advancements,xu2024survey,jodhani2026ai,ejaz2025mlpipeline}. These models are attractive because they can capture nonlinear relationships among weather, topography, vegetation, and human variables while remaining more interpretable and computationally accessible than many deep-learning approaches.

Recent reviews suggest that traditional ML methods remain valuable as robust baselines, especially where interpretability, lower computational cost, and limited training data are important \citep{jain2020review,jodhani2026ai,liu2025advancements,ejaz2025mlpipeline}. However, their performance can deteriorate when spatial heterogeneity is high, when labels are imbalanced, or when the task shifts from susceptibility mapping to real-time fire monitoring or complex spatiotemporal forecasting \citep{caron2025ai}.

\subsection{Deep Learning and Spatiotemporal Prediction}

Deep learning has become increasingly important in wildfire research because of its ability to model high-dimensional spatiotemporal data and learn interactions across imagery, fuels, weather, and terrain. Recent review papers group these approaches into image classification, semantic segmentation, time-series prediction, fire spread forecasting, and broader spatiotemporal risk modeling \citep{xu2024survey,xu2025deeplearning,ejaz2025mlpipeline}. Convolutional neural networks, recurrent neural networks, LSTMs, ConvLSTMs, graph neural networks, transformers, and multimodal architectures have been used for ignition likelihood, fire spread, danger estimation, detection, and burned-area progression \citep{liu2025advancements,xu2025deeplearning}. Earlier deep-learning examples such as FireCast demonstrate the potential of deep models for wildfire-spread prediction, while more recent benchmark datasets support more systematic evaluation \citep{radke2019firecast,huot2022nextday}.

The review literature is broadly consistent in its assessment: deep learning improves predictive flexibility and can achieve strong performance, but operational value depends on interpretability, calibration, transferability, data quality, and computational feasibility \citep{jodhani2026ai,liu2025advancements,caron2025ai}. Data-assimilation studies also indicate a shift toward dynamic model--data fusion. For example, Fire-EnSF proposes an ensemble score filter for wildfire spread data assimilation, integrating numerical fire models with observations to improve real-time spread prediction \citep{shi2025fireensf}. Such work shows that wildfire forecasting is moving beyond static mapping toward dynamic and continuously updated systems.

\subsection{Operational Gaps in Wildfire AI}

Although wildfire AI research is growing rapidly, operational deployment still lags behind research output. Caron et al.\ argue that dataset imbalance, weak benchmarking, unsuitable metrics, computational demands, and limited integration into decision-support workflows remain major barriers \citep{caron2025ai}. Their multi-target wildfire-risk proof of concept further emphasizes that operational wildfire decision-making cannot be reduced to a single risk score: it requires simultaneous consideration of meteorological danger, ignition activity, intervention complexity, and resource mobilization \citep{caron2026multitarget}.

These observations are relevant for the present thesis because they indicate that accurate wildfire prediction is necessary but not sufficient. To support decision-making, models must produce interpretable and scalable outputs that can be linked to consequence analysis, including expected annual losses and broader economic or ecological impacts.

\section{From Wildfire Prediction to Economic Loss Assessment}

\subsection{Suppression Costs and Operational Expenditure}

One strand of recent research examines the relationship between wildfire detection, response timing, and suppression costs. Huang and Wichmann use machine learning and double/debiased machine learning to estimate how reporting delays affect suppression expenditures in Alberta, Canada \citep{huang2024suppression}. Their findings show that faster detection can reduce suppression costs, but that suppression savings alone may not fully justify investments in detection. This is important because the value of early detection also includes avoided property losses, ecological impacts, public-health burdens, and indirect economic disruption.

Ardid et al.\ similarly connect forecast skill with economic value in wildfire potential forecasting, demonstrating that prediction quality should be assessed not only statistically but also in terms of decision value \citep{ardid2026forecastskill}. This perspective is important for a loss-oriented thesis because it shifts evaluation from accuracy alone toward the consequences of improved prediction.

\subsection{Expected Annual Loss and Public Risk Datasets}

A more comprehensive approach explicitly links hazard modeling with exposure and vulnerability to estimate losses. The supranational study by Asif et al.\ develops a machine-learning framework for estimating wildfire losses under climate change in Southeastern Europe \citep{asif2026supranational}. Their method combines random-forest-based susceptibility mapping, hazard classification, exposure layers, vulnerability tables, and economic valuation to estimate average annual loss. This study is especially close to the present thesis because it bridges hazard mapping and loss estimation rather than treating them as separate research streams.

For scalable wildfire-loss research, public risk datasets are also important because they provide standardized targets and exposure-related indicators across large regions. The FEMA National Risk Index provides county- and census-tract-level indicators for multiple natural hazards, including wildfire \citep{fema2025nri_technical,fema2025nri_methodology}. The NRI framework combines expected annual loss, social vulnerability, and community resilience to characterize natural-hazard risk. For wildfire, this makes it possible to move from a purely occurrence-based target toward a monetary risk target that reflects expected consequences.

Expected annual loss is useful because it provides a comparable risk measure across counties and hazards. Instead of asking only whether a wildfire may occur, EAL asks what average annual loss may be expected given hazard frequency, exposure, and consequence assumptions. At the same time, public EAL datasets should be interpreted carefully. They are spatially aggregated, depend on embedded assumptions, and may not capture event-level damage mechanisms directly. They are therefore most appropriate as scalable first baselines that should eventually be complemented by more detailed event-level or asset-level data.

\subsection{Property Damage, Economy-Wide Effects, and Supply Chains}

Beyond suppression and EAL frameworks, wildfire-loss modeling increasingly includes direct structural damage, supply-chain disruption, and economy-wide impacts. Jia and Opabola develop an interpretable machine-learning framework for wildfire damage drivers in California, using explainable AI to identify how hazard, exposure, and vulnerability jointly shape damage outcomes \citep{jia2025interpretable}. Raveendran et al.\ propose a flexible semiparametric approach to wildfire-loss modeling, linking actuarial loss modeling with more flexible statistical methods \citep{raveendran2025wildfireloss}.

Wang et al.\ estimate the economic footprint of the 2018 California wildfires and show that total damages extend far beyond direct fire impacts, including health and indirect economic effects \citep{wang2021economicfootprint}. Monge et al.\ quantify economy-wide wildfire impacts under changing wildfire climate conditions in a regional rural landscape \citep{monge2023probabilistic}. Ma et al.\ extend the loss perspective to supply-chain networks, where wildfire affects physical assets, system performance, transportation times, unmet demand, and total supply-chain costs \citep{ma2022supplychain}. Together, these studies show that wildfire-loss research is moving beyond burned area and direct property damage toward cascading, functional, and economy-wide consequence analysis.

\subsection{Smoke-Related Health and Mortality Burdens}

A rapidly developing component of wildfire-loss literature concerns smoke-related mortality and health damages. Earlier health-impact reviews and global mortality studies show that wildfire smoke exposure can cause substantial health burdens across regions \citep{johnston2012global,reid2016critical}. Burke et al.\ discuss the changing risk and burden of wildfire in the United States, connecting smoke and exposure to broader societal consequences \citep{burke2021smoke}. More recent studies quantify the contribution of anthropogenic climate change to wildfire particulate matter and related mortality, and estimate future wildfire smoke exposure and mortality burdens under climate change \citep{law2025anthropogenic,qiu2025wildfiresmoke}.

This strengthens the argument that wildfire economic losses are not limited to fire-front impacts. They propagate through public health, regional healthcare burdens, long-range smoke exposure, and productivity effects. For a thesis focused on wildfire losses, smoke-health literature broadens the definition of consequence and highlights the need for integrated models that consider both direct and indirect impacts.

\section{Methodological Lessons from Agricultural and Natural-Hazard Loss Modeling}

Agricultural-loss studies provide useful analogies for wildfire-loss research because they explicitly connect environmental hazards with measurable impacts on production, asset value, and recovery. In contrast to many wildfire studies that stop at susceptibility or hazard mapping, this literature often estimates the consequences of hazard exposure. Lateef et al.\ combine optical and SAR imagery with participatory mapping to model flooded crop damage, loss, and recovery, thereby linking physical hazard detection with post-disaster consequence assessment \citep{lateef2025floodedcrop}. Dang et al.\ construct a crop flood damage assessment index for rapid post-disaster impact estimation \citep{dang2024cropflood}.

This direction is reinforced by studies connecting climate extremes to broader agricultural and socioeconomic losses. Zafar et al.\ estimate cropland and infrastructure damage caused by flooding through remote sensing \citep{zafar2024flooddamages}, while Shi et al.\ quantify crop yield and production responses to drought and flood disasters across China \citep{shi2021cropyield}. M\"uller et al.\ show that Sentinel-2 data can quantify drought-related crop-yield impacts at local and regional scales \citep{mueller2025droughtsweden}, and Zhu et al.\ integrate crop, groundwater, and economic dynamics to analyze adaptation and economic outcomes under climate stress \citep{zhu2025ogallala}. Although these studies are centered on agriculture rather than wildfire, they offer transferable methodological lessons: hazard signals must be connected to exposed systems, quantified consequences, and recovery dynamics if risk modeling is to become genuinely loss-oriented.

\section{Environmental Drivers of Wildfire Risk and Loss}

The environmental component of wildfire risk has been extensively documented in climate and fire literature. Variables such as precipitation, vapor pressure deficit, temperature, wind, vegetation structure, fuel type, soil moisture, elevation, and slope influence fuel moisture, fire behavior, ignition potential, spread, and suppression difficulty. These variables are therefore useful not only for predicting fire occurrence but also for loss-oriented wildfire-risk modeling, because they influence where damaging fires are more likely to occur and how severe those events may become.

VPD is particularly important because it captures atmospheric drying demand and is closely related to fuel aridity. Abatzoglou and Williams show that anthropogenic climate change increased fuel aridity and contributed to increased forest fire area across the western United States \citep{abatzoglou2016anthropogenic}. Williams et al.\ further show that atmospheric aridity played a major role in observed increases in California wildfire activity \citep{williams2019observed}. Holden et al.\ add that soil moisture and plant hydraulic stress can be strong predictors of forest fire spread potential, highlighting the importance of vegetation water stress and hydroclimate variables \citep{holden2025soilmoisture}. Beltrán-Marcos et al.\ show that extreme burned area and severity can be driven by atmospheric conditions such as high VPD, strong winds, drought, and heatwave periods \citep{beltranmarcos2026drivers}.

Precipitation is also central because it affects both fuel moisture and vegetation growth. In some settings, low precipitation directly increases dryness, while in others antecedent precipitation increases biomass and later fuel loading. Land-cover variables such as forest, shrub, grass, and developed fractions characterize the exposed landscape and potential fuel structure. Topographic variables such as elevation and slope influence climate gradients and fire behavior, while wind affects spread rate and suppression difficulty. For this reason, geospatial datasets such as gridMET and NLCD are useful for constructing scalable environmental predictors \citep{abatzoglou2013gridmet,dewitz2021nlcd}.

For the present thesis, this literature is directly relevant because the experimental results show that environmental variables substantially improve county-level wildfire EAL prediction. The literature provides scientific explanation for why predictors such as warm-season precipitation, VPD, shrub fraction, elevation, slope, and wind speed are important in the combined model.

\section{Spatial Transferability and Geographic Validation}

A major challenge in wildfire modeling is spatial transferability. Models trained and tested using random splits may appear accurate because nearby or environmentally similar locations can appear in both training and test sets. However, this can overestimate performance when the intended use case involves prediction in new regions or under different environmental conditions. This issue is widely discussed in spatial ecology, environmental modeling, and geospatial machine learning.

Roberts et al.\ argue that blocked or structured cross-validation is often more appropriate than random cross-validation when data have spatial, temporal, hierarchical, or phylogenetic structure \citep{roberts2017crossvalidation}. Valavi et al.\ introduce the blockCV framework for generating spatially or environmentally separated folds in species distribution modeling, emphasizing that random validation can underestimate prediction error in spatially structured data \citep{valavi2019blockcv}. Meyer and Pebesma discuss the problem of predicting into unknown space and propose the concept of area of applicability, where model predictions are more reliable when new data fall within the environmental conditions represented in the training set \citep{meyer2021aoa}. Wang et al.\ further develop spatial+ cross-validation ideas for geospatial problems, reinforcing the need to evaluate models under spatially realistic validation settings \citep{wang2023spatialplus}.

These ideas are important for wildfire-risk modeling because fire regimes, vegetation types, climate gradients, suppression practices, and exposure patterns vary strongly across regions. A model that performs well under random validation may be less reliable when evaluated across held-out states or regions. For this reason, the present thesis uses state-wise grouped validation as a stricter test of geographic generalization. This validation design is not a replacement for all spatial validation methods, but it provides a more demanding benchmark than random splitting and helps assess whether the model captures transferable wildfire-risk signals.

\section{Ecological Losses, Vulnerability, and Recovery}

Compared with economic-loss modeling, ecological-loss assessment remains difficult because ecosystem damage cannot be fully represented by a single monetary measure. Ecological losses may involve biodiversity decline, vegetation degradation, soil degradation, hydrological changes, reduced habitat quality, loss of ecosystem-service capacity, and delayed recovery after severe fire. The European ecological-vulnerability framework proposed by Arrogante-Funes et al.\ is particularly relevant because it models vulnerability through ecological value before fire, coping capacity during fire, and post-fire recovery time \citep{arrogantefunes2024ecological}. Instead of treating burned area as the full proxy for ecological impact, this framework emphasizes that ecological consequences depend on ecosystem value, resistance, coping capacity, and recovery.

\subsection{Ecological Vulnerability and Ecosystem Services}

Recent ecosystem-service studies provide a broader way to interpret ecological wildfire losses. Roces-D\'iaz et al.\ synthesize global evidence on fire effects on ecosystem services in forests and woodlands, showing that fire effects vary across service types, ecosystem contexts, and fire characteristics \citep{rocesdiaz2022ecosystemservices}. Pacheco et al.\ characterize wildfire impacts on ecosystem services in Portugal by triangulating scientific literature, governmental reports, and expert perceptions, demonstrating that wildfire impacts are not limited to vegetation loss but also affect water regulation, erosion control, recreation, cultural values, and provisioning services \citep{pacheco2023ecosystemservices}. Taboada et al.\ show that wildfire can reduce ecosystem-service delivery in maritime pine-dominated forests, including provisioning and cultural services \citep{taboada2021ecosystemdelivery}. Misseyanni et al.\ further review forest-fire impacts on ecosystem services in Greece, emphasizing negative effects on hydrological regulation, erosion control, cultural services, climate regulation, and provisioning services \citep{misseyanni2025ecosystemservices}.

These studies are important for the present thesis because they show that ecological losses are multidimensional. A wildfire may produce direct vegetation damage, but its consequences also extend to ecosystem functions and services that support human and natural systems. Therefore, ecological loss should be treated as a separate but connected component of wildfire-loss assessment rather than as a secondary effect of economic damage.

\subsection{Post-Fire Vegetation Recovery and Remote Sensing}

Post-fire recovery is a central dimension of ecological loss because two areas with similar burned area may differ greatly in recovery speed and long-term ecological degradation. P\'erez-Cabello et al.\ review remote-sensing techniques for assessing post-fire vegetation recovery and show that satellite time series, spectral indices, and multi-temporal analysis can support monitoring of vegetation response after fire \citep{perezcabello2021postfire}. Pe\~na-Molina et al.\ use a 33-year satellite time series to assess fire vulnerability, resilience, and recovery rates of Mediterranean pine forests, showing how long-term Earth-observation data can quantify recovery trajectories \citep{penamolina2024recovery}. Roche et al.\ analyze post-fire recovery dynamics and resilience of ecosystem-service capacity in Mediterranean-type ecosystems, connecting vegetation recovery with ecosystem-service restoration \citep{roche2024postfire}.

This literature is methodologically relevant because it demonstrates how ecological consequences can be measured through remote sensing, time-series analysis, and recovery indicators. For a loss-oriented wildfire thesis, such work suggests that future ecological-loss extensions could combine burn severity, vegetation recovery metrics, land-cover change, and ecosystem-service indicators to complement monetary loss targets.

\subsection{Biodiversity, Multifunctionality, and Long-Term Recovery}

A further ecological dimension concerns biodiversity and ecosystem multifunctionality. Laterza et al.\ show that wildfire severity and time since fire jointly shape biodiversity, ecosystem-service provision, and multifunctionality in Mediterranean forests \citep{laterza2026multifunctionality}. This is important because it indicates that ecological loss depends not only on whether a fire occurred, but also on severity, time since disturbance, and the recovery of multiple ecological functions. Busari et al.\ model wildfire effects on ecosystem services in two California watersheds and communities, illustrating how ecosystem-service impacts can be evaluated spatially after wildfire \citep{busari2025ecosystemservices}. Together, these studies strengthen the argument that ecological consequences require explicit indicators of severity, recovery, and ecosystem services rather than burned area alone.

For the present thesis, this expanded ecological literature suggests that ecological losses should be treated as a distinct but connected dimension of wildfire-loss assessment. Ecological damage may overlap with economic damage, but it is not reducible to it. A complete wildfire-loss framework should therefore incorporate ecological vulnerability, ecosystem services, recovery rates, and landscape-specific resilience alongside monetary impacts.

\section{LLMs, RAG, and Foundation Models in Wildfire Research}

\subsection{LLMs for Literature Synthesis and Decision Support}

Since 2024, wildfire research has begun to incorporate LLMs, multimodal systems, and foundation-model architectures. One use case is large-scale literature synthesis. Lin et al.\ use LLMs to analyze more than 60,000 wildfire research papers and identify geographic and thematic disparities in wildfire research \citep{lin2024llmresearch}. This is relevant because it shows that LLMs can help synthesize large research corpora, although this role is different from quantitative wildfire prediction.

Another use case is decision support. Xie et al.\ develop WildfireGPT, a domain-tailored LLM/RAG system designed to provide context-specific wildfire risk and resilience information \citep{xie2024wildfiregpt,xie2025rag}. Caron et al.\ propose a multi-target wildfire-risk prediction framework in which multiple predictive outputs are synthesized by LLMs into actionable reports for first responders \citep{caron2026multitarget}. Cheng et al.\ show that retrieval-informed prompting can improve agreement between LLM-generated rankings and expert-defined risk rankings in hierarchical wildfire risk evaluation \citep{cheng2026retrieval}. These studies suggest that LLMs are most useful when organizing, retrieving, and communicating complex wildfire knowledge rather than replacing domain-specific quantitative models.

\subsection{Vision-Language and Multimodal Systems}

Multimodal wildfire systems are also becoming more common. Ayanzadeh et al.\ introduce WildfireVLM, which combines satellite-based object detection with language-driven risk assessment and response recommendations \citep{ayanzadeh2026wildfirevlm}. Seidel et al.\ show how vision-language models can enhance UAV-based early wildfire detection by generating structured scene descriptions and improving situational awareness \citep{seidel2025advancing}. FireScope uses chain-of-thought-style conditioning to improve structured wildfire risk prediction and generalization \citep{markov2026firescope}. Xu et al.\ review generative AI for 2D and 3D wildfire spread prediction and argue that generative and multimodal models may become useful complements to physics-based and deep-learning workflows \citep{xu2025generativeai}.

\subsection{Limits of LLMs for Quantitative Wildfire Prediction}

Although these systems are promising, recent studies also make their limitations clear. Ramesh et al.\ compare WildfireGPT-style generative approaches with a domain-specific TabNet model for quantitative fire radiative power forecasting and show that the general-purpose generative system performs poorly for precise numerical prediction \citep{ramesh2025assessingwildfiregpt}. This suggests that LLMs and RAG systems should currently be viewed as augmentation layers for retrieval, explanation, and structured decision support, rather than standalone engines for quantitative wildfire forecasting or loss estimation.

This conclusion is particularly important for the present thesis. LLMs may help translate wildfire-risk and loss information into interpretable reports, but the underlying quantitative backbone still needs to come from domain-specific predictive and loss-estimation models.

\section{Summary of Related Work}

Table~\ref{tab:related_work_summary} summarizes the main categories of related work and their relevance to the present thesis.

\begin{table}[htbp]
\centering
\caption{Summary of related work categories and relevance to this thesis.}
\label{tab:related_work_summary}
\resizebox{\textwidth}{!}{%
\begin{tabular}{p{3.2cm}p{4.2cm}p{4.2cm}p{5.5cm}}
\hline
Category & Representative studies & Main focus & Relevance to thesis \\
\hline
Fire behavior and fire modeling & \citet{rothermel1972spread}, \citet{vanwagner1987fwi}, \citet{finney1998farsite}, \citet{scott2005fuelmodels} & Fuel, weather, topography, fire-spread systems & Provides physical basis for environmental predictors \\
Climate and fire regimes & \citet{westerling2006warming}, \citet{jolly2015climate}, \citet{abatzoglou2016anthropogenic}, \citet{williams2019observed} & Climate, aridity, fire-weather seasons, fire regimes & Justifies VPD, precipitation, temperature, and wind variables \\
Wildfire ML/DL prediction & \citet{jain2020review}, \citet{xu2024survey}, \citet{xu2025deeplearning}, \citet{ejaz2025mlpipeline}, \citet{caron2025ai} & Ignition, susceptibility, detection, spread, prediction pipelines & Provides modeling background and shows limitations of hazard-only prediction \\
Datasets and remote sensing & \citet{eidenshink2007mtbs}, \citet{giglio2016modis}, \citet{schroeder2014viirs}, \citet{xu2024sen2fire}, \citet{karlsson2025modisviirs} & Burn severity, active fire, satellite benchmarks & Supports wildfire data foundations and remote-sensing feature logic \\
Loss and risk assessment & \citet{asif2026supranational}, \citet{fema2025nri_technical}, \citet{huang2024suppression}, \citet{jia2025interpretable}, \citet{raveendran2025wildfireloss} & EAL, suppression costs, damage drivers, loss modeling & Directly supports the thesis target and loss-oriented framing \\
Cascading and indirect impacts & \citet{wang2021economicfootprint}, \citet{monge2023probabilistic}, \citet{ma2022supplychain}, \citet{qiu2025wildfiresmoke} & Economy-wide, supply-chain, smoke-health consequences & Broadens loss definition beyond direct burned area \\
Agricultural-hazard losses & \citet{lateef2025floodedcrop}, \citet{dang2024cropflood}, \citet{shi2021cropyield}, \citet{zhu2025ogallala} & Hazard-to-damage and recovery modeling & Provides methodological analogies for consequence-oriented modeling \\
Ecological losses and recovery & \citet{arrogantefunes2024ecological}, \citet{rocesdiaz2022ecosystemservices}, \citet{perezcabello2021postfire}, \citet{laterza2026multifunctionality} & Ecosystem vulnerability, ecosystem services, vegetation recovery, biodiversity & Strengthens the ecological-loss dimension of the thesis beyond monetary impacts \\
Spatial validation & \citet{roberts2017crossvalidation}, \citet{valavi2019blockcv}, \citet{meyer2021aoa}, \citet{wang2023spatialplus} & Blocked, grouped, and spatial validation & Justifies state-wise grouped validation in the experiments \\
LLMs and foundation models & \citet{xie2024wildfiregpt}, \citet{lin2024llmresearch}, \citet{cheng2026retrieval}, \citet{ramesh2025assessingwildfiregpt} & Retrieval, synthesis, decision support, limitations & Frames LLMs as support tools rather than replacements for quantitative models \\
\hline
\end{tabular}%
}
\end{table}

\section{Research Gap and Thesis Positioning}

The literature demonstrates substantial progress in wildfire occurrence modeling, remote-sensing-based prediction, spatiotemporal deep learning, ecological vulnerability mapping, ecosystem-service impact assessment, public risk indexing, and AI-assisted decision support. However, several gaps remain clear. First, much of the wildfire literature is still hazard-centric, with fewer studies translating predictive outputs into integrated economic and ecological loss estimates \citep{asif2026supranational,caron2025ai}. Second, although public risk datasets such as the FEMA NRI provide useful expected annual loss targets, more work is needed to connect these targets with environmental predictors and machine-learning validation frameworks \citep{fema2025nri_technical,fema2025nri_methodology}. Third, although adjacent natural-hazard fields demonstrate how remote sensing, damage indices, participatory mapping, and economic analysis can estimate consequences and recovery, equivalent integrated pipelines remain less common in wildfire studies \citep{lateef2025floodedcrop,dang2024cropflood,shi2021cropyield,zhu2025ogallala}. Fourth, ecological-loss studies increasingly examine ecosystem services, post-fire recovery, biodiversity, and multifunctionality, but these indicators are still rarely integrated with quantitative wildfire-risk and monetary-loss models \citep{rocesdiaz2022ecosystemservices,perezcabello2021postfire,roche2024postfire,laterza2026multifunctionality}. Fifth, spatial transferability remains an important open challenge because wildfire models can perform well under random validation but degrade when evaluated across new geographic regions \citep{roberts2017crossvalidation,valavi2019blockcv,meyer2021aoa,wang2023spatialplus}. Finally, while LLMs and foundation models are becoming useful for retrieval, synthesis, and interpretation, they are not yet reliable substitutes for quantitative wildfire-risk and wildfire-loss models \citep{xie2025rag,ramesh2025assessingwildfiregpt,cheng2026retrieval}.

Accordingly, the present thesis is positioned at the intersection of wildfire-risk prediction and wildfire-loss assessment. Its contribution lies not in treating wildfire hazard as an endpoint, but in viewing it as an input to a broader framework that considers what is exposed, how vulnerable it is, what environmental conditions shape wildfire likelihood and severity, and what economic and ecological consequences follow. The thesis responds to the need for integrated, geographically validated, consequence-oriented wildfire-loss modeling using accessible public datasets.

\chapter{Methodology}
\label{chap:methodology}

\section{Overview}

This chapter describes the methodology used to construct and evaluate the county-level wildfire-loss modeling pipeline. The goal of the methodology is to move from a hazard-oriented framing of wildfire prediction toward a loss-oriented setup in which the target variable is county-level wildfire expected annual loss (EAL). The workflow combines publicly accessible tabular data, geospatial environmental predictors, and machine-learning regression models.

The methodology was implemented in a sequence of stages. First, a cleaned county-level wildfire-loss target table was constructed from the FEMA National Risk Index (NRI), which provides standardized natural-hazard risk indicators including expected annual loss, social vulnerability, and community resilience \citep{fema2025nri_technical,fema2025nri_methodology}. Second, a baseline predictor dataset was built using county-level demographic, exposure, and resilience variables. Third, the baseline was extended with wildfire-relevant environmental predictors derived from public geospatial datasets. Fourth, several regression models were evaluated under both conventional random train/test splits and stricter state-wise grouped validation.

\section{Study Design}

The current experimental design is county-based. Each observation corresponds to one county or county-equivalent unit, and the objective is to predict wildfire expected annual loss in monetary units. This design was chosen for three reasons. First, county-level data are broadly accessible and allow reproducible modeling at national scale. Second, county-level units make it possible to merge monetary loss targets with exposure, vulnerability, and environmental predictors. Third, the county scale provides a practical intermediate level between highly aggregated state analysis and much more detailed event-level or parcel-level loss modeling.


It is important to clarify the empirical scope of the methodology. The implemented modeling pipeline focuses on monetary wildfire Expected Annual Loss as the prediction target. Ecological losses are not directly modeled as a separate dependent variable because the current FEMA NRI-based dataset does not provide a county-level ecological-loss target comparable to monetary EAL. However, ecological context is partially represented through environmental predictors such as land-cover fractions, elevation, slope, precipitation, wind speed, and vapor pressure deficit. These variables help characterize wildfire hazard conditions, but they should not be interpreted as direct ecological-loss measurements.

The present modeling design should be viewed as a county-level loss-assessment framework rather than an event-level or asset-level damage model. It is intended to test whether wildfire expected annual loss can be explained to a meaningful extent with accessible public data and whether environmental predictors improve performance beyond purely county-level socioeconomic variables.

\section{Complete Pipeline for Wildfire Risk Prediction and Assessment}
\label{sec:complete_pipeline}

The proposed methodology is designed as an end-to-end pipeline for county-level wildfire risk prediction and assessment. The purpose of the pipeline is not only to train a regression model, but to convert heterogeneous public datasets into a structured wildfire-risk assessment output in monetary units. The complete workflow connects data acquisition, target construction, feature engineering, environmental augmentation, model training, validation, prediction, and risk assessment.

The pipeline begins with public county-level risk data from the FEMA National Risk Index. From this source, wildfire Expected Annual Loss (EAL) is extracted as the main target variable. This target is suitable for the thesis because it represents wildfire risk in monetary units rather than only fire occurrence or burned area. The raw wildfire EAL values are cleaned, standardized by county FIPS identifiers, and transformed using a logarithmic transformation to reduce skewness.

The second stage constructs the baseline county-level predictor table. This includes demographic, exposure, value, vulnerability, and resilience-related variables such as population, building value, agriculture value, county area, social vulnerability, and community resilience. Derived density-style variables are then created in order to represent value concentration and exposure intensity at county level.

The third stage adds environmental predictors. Since wildfire losses depend not only on exposed assets but also on environmental conditions that influence wildfire occurrence and severity, the baseline table is augmented with topographic, land-cover, and climate variables. These include elevation, slope, forest fraction, shrub fraction, grass fraction, developed-land fraction, warm-season maximum temperature, precipitation, wind speed, and vapor pressure deficit. These predictors represent key components of the wildfire hazard environment.

The fourth stage trains machine-learning regression models to predict the transformed wildfire EAL target. Linear Regression and Ridge Regression are used as simple baseline models, while Random Forest and Gradient Boosting are used to capture nonlinear relationships among exposure, vulnerability, resilience, and environmental variables. The best-performing model is then used as the main predictive component of the pipeline.

The fifth stage evaluates the model using both random and geographically stricter validation strategies. Random train/test splitting and standard cross-validation provide an initial performance benchmark. State-wise grouped validation is then used to test whether the model can generalize across geographic regions. This is important because a wildfire-risk model should not only interpolate among nearby counties but should also retain predictive value when applied to held-out states.

The final stage converts model predictions into risk-assessment outputs. The model predicts log-transformed wildfire EAL, which can be converted back into monetary units using the inverse transformation:

\begin{equation}
\widehat{\mathrm{wildfire\_eal}} = \exp(\hat{y}) - 1.
\end{equation}

These predicted EAL values can then be used to rank counties, compare observed and predicted losses, identify counties with high expected wildfire losses, and support broader risk interpretation. In this sense, the pipeline produces not only a predictive model but also a structured risk-assessment framework.

Overall, the complete pipeline can be summarized as:
\[
\begin{aligned}
&\text{Public data sources}
\rightarrow
\text{cleaned county table}
\rightarrow
\text{baseline predictors} \\
&\rightarrow
\text{environmental augmentation}
\rightarrow
\text{ML regression model} \\
&\rightarrow
\text{geographic validation}
\rightarrow
\text{wildfire risk assessment output}.
\end{aligned}
\]

Figure~\ref{fig:fema_nri_risk_pipeline} summarizes the complete FEMA/NRI-based workflow used in this thesis. The diagram shows how public risk, exposure, vulnerability, resilience, topographic, land-cover, and climate datasets are transformed into predicted wildfire EAL and county-level risk-assessment outputs.

\begin{figure}[htbp]
    \centering
    \safefigure[width=\textwidth]{figures/pipelines/wildfire_loss_assessment_workflow_diagram.pdf}
    \caption{Complete FEMA/NRI-based pipeline for county-level wildfire risk prediction and assessment. The workflow transforms public risk, exposure, vulnerability, resilience, topographic, land-cover, and climate data into predicted wildfire Expected Annual Loss and risk-assessment outputs.}
    \label{fig:fema_nri_risk_pipeline}
\end{figure}

Before the current county-level FEMA NRI pipeline was developed, an exploratory fine-scale Creek Fire data-engineering workflow was prepared as a baseline approach for wildfire tensor construction. Since the final thesis direction focuses on wildfire risk/loss assessment in monetary units, the Creek Fire workflow is included in Appendix~\ref{app:creek_fire_pipeline} as supporting preliminary work rather than as the main experimental pipeline.

\section{Target Variable Construction}

\subsection{Wildfire Expected Annual Loss}

The main target variable is county-level wildfire expected annual loss, denoted as \texttt{wildfire\_eal}. This variable was extracted from the FEMA National Risk Index county table. The NRI was selected because it provides a publicly accessible consequence-oriented target for large-scale modeling and is structured around expected annual loss, social vulnerability, and community resilience \citep{fema2025nri_technical,fema2025nri_methodology}. This makes it suitable for a scalable experiment in loss-oriented wildfire-risk modeling.

It is important to clarify that FEMA NRI wildfire EAL is not a raw record of observed wildfire damages from individual fire events. It is a modeled expected-loss indicator that reflects the combination of hazard frequency, exposure, and expected consequence. Therefore, the present thesis does not predict individual wildfire-event damages directly. Instead, it learns county-level relationships between publicly available exposure, vulnerability, resilience, environmental predictors, and the FEMA NRI wildfire EAL target. This framing makes the target suitable for county-level risk assessment, but it also means that the results should be interpreted as expected-loss modeling rather than event-level damage reconstruction.

The raw target distribution is strongly right-skewed, with many counties having relatively low wildfire EAL values and a smaller number of counties having substantially larger losses. To stabilize the regression problem, the target was transformed using
\begin{equation}
y = \log(1 + \texttt{wildfire\_eal}).
\end{equation}
This transformed target, denoted as \texttt{log\_wildfire\_eal}, was used in all machine-learning experiments.

\subsection{Initial Cleaning}

The original county-level table contained 3232 geographic units. Rows with missing wildfire target values were removed, producing a cleaned target dataset with 3144 rows. County identifiers were standardized using 5-digit FIPS codes in string format in order to ensure consistent merging across all later datasets.

Table~\ref{tab:nri_target_construction_summary_export} summarizes the target construction process. It reports the original number of FEMA NRI county records, the number of rows removed because of missing wildfire EAL values, and the final number of counties retained for modeling.

\begin{table}[htbp]
    \centering
    \safetablepdf[width=0.9\textwidth]{tables/notebook1/table_1_target_construction_summary.pdf}
    \caption{Target construction summary from the FEMA NRI county-level wildfire EAL table.}
    \label{tab:nri_target_construction_summary_export}
\end{table}

\section{Baseline County-Level Predictors}

The first predictor set was designed as a county-level baseline using accessible variables already present in the county table. These variables were selected to represent exposure, value, vulnerability, resilience, and spatial scale.

Table~\ref{tab:selected_baseline_predictors_export} lists the selected raw FEMA/NRI fields, their corresponding model variable names, their role in the thesis, and the reason for including each variable. The selected baseline predictors include population, building value, agriculture value, county area, social vulnerability score, and community resilience score.

\begin{table}[htbp]
    \centering
    \caption{Selected baseline county-level predictors used to represent exposure, value, vulnerability, and resilience.}
    \label{tab:selected_baseline_predictors_export}
    \scriptsize
    \setlength{\tabcolsep}{4pt}
    \begin{tabularx}{\textwidth}{>{\raggedright\arraybackslash}p{2.3cm} >{\raggedright\arraybackslash}p{3.0cm} >{\raggedright\arraybackslash}p{2.6cm} X}
    \toprule
    Raw FEMA/NRI field & Model variable & Role in thesis & Reason for inclusion \\
    \midrule
    \texttt{POPULATION} & \texttt{population} & Population exposure & Represents exposed population at county level. \\
    \texttt{BUILDVALUE} & \texttt{building\_value} & Building exposure / asset value & Represents built-environment value exposed to wildfire risk. \\
    \texttt{AGRIVALUE} & \texttt{agriculture\_value} & Agricultural exposure / asset value & Represents agricultural value exposed to wildfire risk. \\
    \texttt{AREA} & \texttt{area\_sq\_mi} & County spatial scale & Controls for county size and spatial scale. \\
    \texttt{SOVI\_SCORE} & \texttt{social\_vulnerability\_
    score} & Social vulnerability & Captures the social vulnerability component of risk. \\
    
    \texttt{RESL\_SCORE} & \texttt{community\_
    resilience\_score} & Community resilience & Captures community capacity and resilience. \\
    \bottomrule
    \end{tabularx}
\end{table}

In addition to these direct variables, several derived features were constructed:
\begin{itemize}
    \item population density,
    \item building value density,
    \item agriculture value density,
    \item building value per capita.
\end{itemize}

Table~\ref{tab:derived_baseline_features_export} summarizes these derived variables, their formulas, and their interpretation. These derived features provide an interpretable approximation of exposure intensity and value concentration, which are important for loss-oriented wildfire modeling.

\begin{table}[htbp]
    \centering
    \safetablepdf[width=\textwidth]{tables/notebook2/table_2_derived_baseline_features.pdf}
    \caption{Derived baseline features used to represent exposure intensity and value concentration.}
    \label{tab:derived_baseline_features_export}
\end{table}

After merging the target and baseline predictors, leakage-prone wildfire-specific risk fields were removed from the modeling table. The excluded fields were:

\begin{itemize}
    \setlength{\itemsep}{0pt}
    \item \texttt{wildfire\_risk\_score}
    \item \texttt{wildfire\_risk\_rating}
    \item \texttt{wildfire\_eal\_score\_like}
    \item \texttt{wildfire\_eal\_rating\_like}
    \item \texttt{wildfire\_risk\_value}
    \item \texttt{wildfire\_eal\_buildings}
    \item \texttt{wildfire\_eal\_population}
    \item \texttt{wildfire\_eal\_agriculture}
    \item \texttt{wildfire\_annual\_frequency}
\end{itemize}

These variables were excluded because they are directly derived from, or closely related to, the FEMA wildfire-risk calculation. Keeping them would allow the model to learn from fields that already encode wildfire-specific risk or loss information, making the prediction task artificially easier. Table~\ref{tab:final_baseline_dataset_summary_export} reports the final non-leaky baseline modeling dataset summary after target--predictor merging.

\begin{table}[htbp]
    \centering
    \safetablepdf[width=0.85\textwidth]{tables/notebook2/table_3_final_baseline_dataset_summary.pdf}
    \caption{Summary of the final non-leaky baseline modeling dataset after target--predictor merging.}
    \label{tab:final_baseline_dataset_summary_export}
\end{table}

\section{Environmental Augmentation}

\subsection{Motivation}

The baseline county variables describe exposure and resilience, but they do not fully capture wildfire-relevant environmental conditions. Since wildfire activity and its consequences depend strongly on terrain, fuels, vegetation structure, climate, atmospheric dryness, soil moisture, and wind, the next stage of the methodology added environmental predictors derived from public geospatial datasets. This step was also motivated by wildfire-climate and fire-environment literature showing the importance of variables such as vapor pressure deficit, precipitation, fuel aridity, soil moisture, plant water stress, wind, and drought conditions for wildfire activity and spread \citep{abatzoglou2016anthropogenic,williams2019observed,holden2025soilmoisture,beltranmarcos2026drivers}.

\subsection{Spatial Scope}

Environmental augmentation was initially performed for the contiguous United States (CONUS). This restriction was introduced because the selected climate and land-cover datasets were directly compatible with a CONUS-focused workflow. The resulting environmental merge excluded a small number of unmatched Connecticut county-equivalent units, which were documented separately.

\subsection{Environmental Data Sources}

Environmental predictors were derived from the following public datasets:
\begin{itemize}
    \item topography from SRTM,
    \item county geometries from TIGER counties,
    \item land cover from the 2019 NLCD release,
    \item warm-season climate summaries from gridMET.
\end{itemize}

Warm-season climate variables were derived from gridMET, a gridded surface meteorological dataset developed for ecological and environmental modeling \citep{abatzoglou2013gridmet}. Land-cover predictors were derived from the National Land Cover Database (NLCD), which provides Landsat-based land-cover products at 30 m spatial resolution \citep{dewitz2021nlcd}.

The selected environmental variables were:
\begin{itemize}
    \item mean elevation,
    \item mean slope,
    \item forest fraction,
    \item shrub fraction,
    \item grass fraction,
    \item developed-land fraction,
    \item warm-season mean maximum temperature,
    \item warm-season mean precipitation,
    \item warm-season mean wind speed,
    \item warm-season mean vapor pressure deficit.
\end{itemize}

Table~\ref{tab:environmental_predictor_definitions_export} defines the environmental predictors used for the CONUS county-level wildfire EAL modeling. These predictors were chosen because they describe key environmental controls relevant to wildfire hazard and consequence patterns, including terrain, vegetation composition, moisture availability, atmospheric dryness, and wind.

% THIS IS MADE THROUGH LATEX
%\begin{table}[htbp]
    %\centering
    %\safetablepdf[width=\textwidth]{tables/notebook5/table_2_environmental_predictor_definitions.pdf}
    %\caption{Environmental predictors used for CONUS county-level wildfire EAL modeling.}
    %\label{tab:environmental_predictor_definitions_export}
%\end{table}

\begin{table}[htbp]
\centering
\caption{Environmental predictors used for CONUS county-level wildfire EAL modeling.}
\label{tab:environmental_predictor_definitions_export}
\scriptsize
\setlength{\tabcolsep}{4pt}
\renewcommand{\arraystretch}{1.15}
\begin{tabularx}{\textwidth}{L{2.6cm} L{2.8cm} L{3.9cm} X}
\toprule
Predictor & Data source & County-level aggregation & Role in wildfire-risk modeling \\
\midrule
\texttt{elevation\_m} & SRTM & Mean over county polygon & Topographic/climatic gradient \\
\texttt{slope\_deg} & SRTM-derived terrain & Mean over county polygon & Terrain influence on fire behavior \\
\texttt{forest\_frac} & NLCD 2019 & Mean binary land-cover fraction & Forest fuel context \\
\texttt{shrub\_frac} & NLCD 2019 & Mean binary land-cover fraction & Shrub fuel context \\
\texttt{grass\_frac} & NLCD 2019 & Mean binary land-cover fraction & Grass/agricultural fuel context \\
\texttt{developed\_frac} & NLCD 2019 & Mean binary land-cover fraction & Human/developed-land context \\
\texttt{warm\_tmmx\_c} & gridMET & Warm-season mean, then county mean & Thermal fire-weather condition \\
\texttt{warm\_pr\_mm} & gridMET & Warm-season mean, then county mean & Moisture/fuel-availability condition \\
\texttt{warm\_wind\_ms} & gridMET & Warm-season mean, then county mean & Potential fire-spread/suppression condition \\
\texttt{warm\_vpd\_kpa} & gridMET & Warm-season mean, then county mean & Atmospheric dryness/fuel-aridity condition \\
\bottomrule
\end{tabularx}
\end{table}

\subsection{County-Level Aggregation}

Environmental layers were aggregated to the county level by spatial reduction over county polygons. For topographic variables such as elevation and slope, county-level mean values were computed. For categorical land-cover information, land-cover classes were converted into fractional predictors such as forest fraction, shrub fraction, grass fraction, and developed-land fraction. For climate variables, warm-season means were computed from daily gridMET observations over a multi-year period and then spatially reduced by county.

NoData or missing raster values were excluded from the aggregation where applicable. The purpose of this step was not to preserve fine-scale pixel variability, but to create consistent county-level environmental predictors that could be merged with the FEMA NRI county-level target using FIPS identifiers. The resulting county-level environmental table was merged with the cleaned county baseline using county FIPS codes.

\subsection{Exploratory Ecological Risk-Screening Layer}

Because the main FEMA/NRI target represents monetary wildfire Expected Annual Loss, the thesis does not directly estimate ecological loss as a separate supervised-learning target. However, an exploratory ecological risk-screening layer was added to connect the monetary risk model with the ecological dimension of wildfire consequences. This layer uses the land-cover fractions already derived during environmental augmentation to represent natural vegetation exposure.

First, a natural vegetation fraction was calculated by combining forest, shrub, and grass fractions. This variable identifies counties where a larger share of the landscape consists of vegetation types that may be ecologically affected by wildfire. Second, the natural vegetation fraction was combined with the predicted wildfire-risk percentile from the final FEMA/NRI model to produce an exploratory ecological risk proxy. The proxy is defined as:
\[
\text{Ecological risk proxy} = \text{Natural vegetation fraction} \times \text{Predicted wildfire-risk percentile}.
\]

This proxy does not represent measured ecological damage, biodiversity loss, burn severity, or ecosystem-service loss. Instead, it identifies counties where high predicted monetary wildfire risk overlaps with high natural vegetation exposure. The proxy therefore provides an ecological screening perspective that complements the main monetary EAL model without changing the primary supervised-learning target.

A second exploratory step classified counties into economic--ecological overlap groups. Counties were categorized according to whether they had high predicted economic wildfire risk and high natural vegetation exposure. This analysis highlights counties where monetary wildfire risk and ecological exposure may coincide, and it separates them from counties where only one of these dimensions is high.

\section{Model Development}

\subsection{Regression Models}

Several regression models were evaluated in the experiments:
\begin{itemize}
    \item Linear Regression,
    \item Ridge Regression,
    \item Random Forest Regression,
    \item Gradient Boosting Regression.
\end{itemize}

The linear models were included as simple baselines, while the tree-based ensemble models were used to capture nonlinear relationships between wildfire EAL and the predictor variables. The experiments showed that Random Forest consistently produced the strongest performance, especially after environmental augmentation.

\subsection{Preprocessing}

The regression pipeline used standard machine-learning preprocessing steps. Numeric predictors were median-imputed when necessary and then standardized before model fitting. Although tree-based models are less sensitive to scaling than linear models, using a consistent preprocessing pipeline allowed direct comparison across multiple algorithms.

\section{Validation Strategy}

\subsection{Random Train/Test Split}

The first validation setup used a conventional random train/test split, with 80\% of observations assigned to training and 20\% assigned to testing. This provided a first estimate of predictive performance under standard machine-learning practice.

\subsection{Cross-Validation}

For baseline benchmarking, 5-fold cross-validation was used to evaluate the stability of the county-only and environmental-augmented models. Performance was summarized through the mean and standard deviation of the fold-wise results.

\subsection{State-Wise Grouped Validation}

Because random splitting can still place geographically similar counties in both training and testing subsets, a stricter geographic validation strategy was also introduced. In this setup, grouped cross-validation by state was used, so that counties from held-out states were not observed during training. This allowed a more realistic test of geographic transferability.

This stricter validation design is motivated by the spatial machine-learning literature, which shows that random validation can overestimate performance when observations are spatially structured or geographically dependent \citep{roberts2017crossvalidation,valavi2019blockcv}. It is also consistent with the broader concern that spatial prediction models may be less reliable when applied outside the environmental conditions represented in the training data \citep{meyer2021aoa}. Recent geospatial validation work further emphasizes that spatially realistic validation strategies are necessary when models are expected to generalize across space rather than only interpolate among nearby observations \citep{wang2023spatialplus}.

\section{Evaluation Metrics}

The main evaluation metrics used in the experiments were:
\begin{itemize}
    \item coefficient of determination ($R^2$),
    \item root mean squared error (RMSE),
    \item mean absolute error (MAE).
\end{itemize}

The coefficient of determination was used as the primary metric for model comparison because it provides an interpretable measure of the fraction of variance explained in the transformed wildfire EAL target. RMSE and MAE were also reported in order to quantify average prediction error and sensitivity to larger residuals.

\section{Ablation and Interpretation}

To better understand how the county-level predictors contributed to performance, an ablation analysis was conducted on the baseline Random Forest model. The ablation tested the effect of removing county area, social vulnerability, and community resilience from the feature set. This allowed the experiments to assess whether predictive performance depended mainly on broad spatial scale, or whether vulnerability and resilience variables added meaningful explanatory value.

Feature importance from the Random Forest models was also used for interpretation. This provided an initial indication of which variables contributed most strongly to the predictive signal in both the county-only and combined county-plus-environment settings.

\section{Methodological Summary}

In summary, the methodology combines three components:
\begin{enumerate}
    \item a county-level wildfire expected annual loss target in monetary units,
    \item a baseline set of county exposure, value, and resilience predictors,
    \item an environmental augmentation layer built from topography, land cover, and climate.
\end{enumerate}

This methodology was designed to test whether publicly accessible datasets are sufficient to support a meaningful first benchmark for wildfire-loss modeling. The resulting experimental findings are presented in Chapter~\ref{chap:results}.

\chapter{Experimental Results and Risk Assessment}
\label{chap:results}

\section{Overview}

This chapter presents the experimental results of the county-level wildfire-loss modeling pipeline in monetary units. Following the literature gap identified in Chapter~\ref{chap:litreview}, the experimental pipeline was designed to move beyond wildfire hazard prediction alone and toward a consequence-oriented formulation in which the prediction target is wildfire \emph{Expected Annual Loss} (EAL). The current experiments use publicly accessible county-level data and therefore serve as a first scalable baseline rather than a final loss model.

The modeling workflow was developed in six connected stages. First, a county-level wildfire EAL target was constructed from the FEMA National Risk Index (NRI). Second, baseline county-level exposure, value, vulnerability, and resilience predictors were prepared. Third, baseline regression models were trained and evaluated. Fourth, ablation and error analyses were conducted to understand the baseline model. Fifth, environmental predictors derived from topography, land cover, and warm-season climate summaries were added. Sixth, the combined model was evaluated under a stricter state-wise grouped validation protocol in order to assess geographic robustness.

\section{Target Variable and Dataset Construction}

The main prediction target in the current experiments is county-level wildfire expected annual loss, denoted as \texttt{wildfire\_eal}, obtained from the FEMA National Risk Index. The NRI provides standardized risk indicators based on expected annual loss, social vulnerability, and community resilience, making it suitable for a first scalable loss-oriented wildfire-risk modeling baseline \citep{fema2025nri_technical,fema2025nri_methodology}. Because the raw EAL values are strongly right-skewed, the modeling target was transformed using
\[
y = \log(1 + \texttt{wildfire\_eal}),
\]
which produced a more stable target distribution for regression.

The initial NRI-based county dataset contained 3232 geographic units. After removing 88 rows with missing wildfire target values, the cleaned target table contained 3144 counties. The exported target construction summary is included in Chapter~\ref{chap:methodology}. Figure~\ref{fig:eal_raw_log_distribution} shows the distributional motivation for applying the logarithmic transformation. The raw wildfire EAL distribution is highly concentrated near low values with a long right tail, whereas the transformed target is more suitable for regression modeling.

\begin{figure}[htbp]
    \centering
    \safefigure[width=\textwidth]{figures/notebook1/figure_1_wildfire_eal_raw_vs_log_distribution.pdf}
    \caption{Distribution of FEMA NRI wildfire Expected Annual Loss before and after logarithmic transformation. The raw target is highly right-skewed, while the log-transformed target is more suitable for regression modeling.}
    \label{fig:eal_raw_log_distribution}
\end{figure}

\section{Baseline County-Level Risk Model}

The first modeling experiment evaluated whether county-level wildfire EAL could be predicted using only accessible county variables without additional environmental augmentation. Several regression models were tested, including Linear Regression, Ridge Regression, Random Forest, and Gradient Boosting. Among these, Random Forest produced the strongest baseline performance.

Table~\ref{tab:baseline_model_performance_summary_export} reports the combined test-set and cross-validation performance of the baseline county-level models. For the baseline county-only model, the Random Forest achieved a test-set $R^2$ of approximately 0.483, with RMSE of approximately 1.675 and MAE of approximately 1.292. Under 5-fold cross-validation, the same model achieved a mean $R^2$ of approximately 0.504 with standard deviation 0.017, and a mean RMSE of approximately 1.661 with standard deviation 0.028. These results indicate that accessible county-level exposure, value, and resilience variables already contain a meaningful signal for predicting wildfire EAL, but they also suggest that a substantial part of the variation remains unexplained.

\begin{table}[htbp]
    \centering
    \safetablepdf[width=\textwidth]{tables/notebook3/table_3_combined_baseline_model_performance_summary.pdf}
    \caption{Combined test-set and cross-validation performance of the baseline county-level wildfire EAL models.}
    \label{tab:baseline_model_performance_summary_export}
\end{table}

Figure~\ref{fig:baseline_cv_r2_comparison} compares the cross-validation $R^2$ values across the baseline models. The Random Forest model achieved the highest mean cross-validation $R^2$, followed closely by Gradient Boosting. Figure~\ref{fig:baseline_cv_rmse_comparison} shows the corresponding cross-validation RMSE comparison, again indicating that the tree-based ensemble models performed better than the linear baselines.

\begin{figure}[htbp]
    \centering
    \safefigure[width=0.82\textwidth]{figures/notebook3/figure_3_cv_r2_baseline_model_comparison.pdf}
    \caption{Cross-validation $R^2$ comparison for baseline county-level models. Random Forest achieved the strongest baseline performance.}
    \label{fig:baseline_cv_r2_comparison}
\end{figure}

\begin{figure}[htbp]
    \centering
    \safefigure[width=0.82\textwidth]{figures/notebook3/figure_4_cv_rmse_baseline_model_comparison.pdf}
    \caption{Cross-validation RMSE comparison for baseline county-level models.}
    \label{fig:baseline_cv_rmse_comparison}
\end{figure}

The weaker performance of Linear Regression and Ridge Regression compared with Random Forest suggests that the relationship between county predictors and wildfire loss is nonlinear. In other words, county-level wildfire risk in monetary units cannot be adequately represented by a simple linear combination of exposure and vulnerability indicators.

\section{Baseline Diagnostics and Feature Importance}

The Random Forest baseline was further analyzed using predicted-versus-actual plots, residual diagnostics, and feature importance. These diagnostics help show whether the baseline model is merely memorizing scale effects or whether it captures meaningful signals related to exposure, vulnerability, and resilience.

Figure~\ref{fig:baseline_predicted_vs_actual} shows the predicted versus actual log-transformed wildfire EAL values for the best baseline model. The scatter pattern indicates that the model captures part of the target variation, but the spread around the diagonal line also confirms that the county-only baseline has limited explanatory power.

\begin{figure}[htbp]
    \centering
    \safefigure[width=0.72\textwidth]{figures/notebook3/figure_5_predicted_vs_actual_best_model.pdf}
    \caption{Predicted versus actual log-transformed wildfire EAL for the best baseline model.}
    \label{fig:baseline_predicted_vs_actual}
\end{figure}

Figure~\ref{fig:baseline_rf_feature_importance} presents the top Random Forest feature importances for the baseline county-level model. The feature-importance results help identify which exposure, value, vulnerability, and resilience variables contributed most strongly before environmental predictors were added.

\begin{figure}[htbp]
    \centering
    \safefigure[width=0.78\textwidth]{figures/notebook3/figure_7_random_forest_feature_importance_top10.pdf}
    \caption{Top Random Forest feature importances for the baseline county-level model.}
    \label{fig:baseline_rf_feature_importance}
\end{figure}

The baseline diagnostics are useful, but they also show the limits of county-only modeling. The baseline model captures meaningful variation but does not yet fully represent environmental wildfire hazard context. This motivated the environmental augmentation step described later in this chapter.

\section{Ablation Analysis of County-Level Variables}

To better understand which county-level variables contributed most to performance, an ablation study was conducted on the baseline Random Forest model. Table~\ref{tab:rf_ablation_performance_export} reports the Random Forest ablation results for different feature-removal settings. The full feature set performed best, with cross-validated $R^2 \approx 0.504$. Removing county area caused only a modest reduction in performance, indicating that although county size contributes to predictive quality, it is not the sole driver of model performance. Removing social vulnerability caused a somewhat larger reduction, while removing community resilience produced a noticeably stronger drop. When both vulnerability and resilience variables were removed together, performance declined to approximately 0.419.

% THIS IS MADE THROUGH LATEX
%\begin{table}[htbp]
    %\centering
    %\safetablepdf[width=\textwidth]{tables/notebook4/table_1_random_forest_ablation_performance.pdf}
    %\caption{Random Forest ablation results for the baseline county-level feature set.}
    %\label{tab:rf_ablation_performance_export}
%\end{table}

\begin{table}[htbp]
\centering
\caption{Random Forest ablation results for the baseline county-level feature set.}
\label{tab:rf_ablation_performance_export}
\scriptsize
\setlength{\tabcolsep}{4pt}
\renewcommand{\arraystretch}{1.12}
\resizebox{\textwidth}{!}{%
\begin{tabular}{lrrrrrrrr}
\toprule
Feature set & No. features & MAE & RMSE & Test $R^2$ & CV $R^2$ mean & CV $R^2$ std & CV RMSE mean & CV RMSE std \\
\midrule
\texttt{all\_features} & 10.0000 & 1.2920 & 1.6748 & 0.4832 & 0.5039 & 0.0173 & 1.6613 & 0.0283 \\
\texttt{without\_area} & 9.0000 & 1.3118 & 1.6745 & 0.4834 & 0.4926 & 0.0160 & 1.6802 & 0.0313 \\
\texttt{simple\_core} & 6.0000 & 1.2965 & 1.6896 & 0.4740 & 0.4891 & 0.0101 & 1.6860 & 0.0230 \\
\texttt{without\_vulnerability} & 9.0000 & 1.3234 & 1.7051 & 0.4643 & 0.4854 & 0.0177 & 1.6920 & 0.0365 \\
\texttt{without\_resilience} & 9.0000 & 1.4144 & 1.8248 & 0.3865 & 0.4338 & 0.0264 & 1.7744 & 0.0349 \\
\texttt{without\_sovi\_resl} & 8.0000 & 1.4475 & 1.8565 & 0.3650 & 0.4192 & 0.0290 & 1.7972 & 0.0432 \\
\bottomrule
\end{tabular}%
}
\end{table}

Figure~\ref{fig:ablation_cv_r2} visualizes the cross-validated $R^2$ values under the ablation settings. The figure makes the performance decline clearer and shows that vulnerability and resilience-related variables add useful predictive information beyond general exposure and spatial-scale variables.

\begin{figure}[htbp]
    \centering
    \safefigure[width=0.82\textwidth]{figures/notebook4/figure_1_ablation_cv_r2.pdf}
    \caption{Cross-validated $R^2$ under different feature-ablation settings.}
    \label{fig:ablation_cv_r2}
\end{figure}

These ablation results are important because they show that the model is not dominated only by generic spatial scale effects. Rather, county-level social and resilience-related variables provide additional signal for wildfire EAL prediction. This supports the broader thesis argument that wildfire losses depend not only on hazard and exposure, but also on underlying vulnerability and resilience conditions.

\section{Environmental Augmentation of the County Model}

The next experiment extended the county baseline by adding wildfire-relevant environmental predictors derived from public geospatial datasets. These included topographic variables, land-cover fractions, and warm-season climate summaries. The goal of this step was to test whether hazard-relevant environmental context could improve the loss-oriented county model.

Table~\ref{tab:environmental_model_comparison_export} reports the model comparison after adding environmental predictors. The results show a substantial performance improvement, especially for the combined baseline-plus-environment feature set. This indicates that the model benefits from combining exposure and vulnerability information with environmental wildfire-hazard context.

\begin{table}[htbp]
    \centering
    \caption{Model comparison after adding environmental predictors to the county-level baseline dataset.}
    \label{tab:environmental_model_comparison_export}
    \scriptsize
    \setlength{\tabcolsep}{4pt}
    \begin{tabularx}{\textwidth}{>{\raggedright\arraybackslash}p{3.1cm} >{\raggedright\arraybackslash}p{2.0cm} *{5}{>{\centering\arraybackslash}p{1.2cm}}}
        \toprule
        Feature set & Model & MAE & RMSE & $R^2$ & CV $R^2$ & CV RMSE \\
        \midrule
        Combined baseline + environment & Random Forest & 0.779 & 1.011 & 0.829 & $0.795 \pm 0.022$ & $1.049 \pm 0.023$ \\
        Combined baseline + environment & Gradient Boosting & 0.864 & 1.085 & 0.804 & $0.776 \pm 0.019$ & $1.099 \pm 0.018$ \\
        Environmental only & Random Forest & 0.817 & 1.112 & 0.794 & $0.769 \pm 0.017$ & $1.118 \pm 0.030$ \\
        Environmental only & Gradient Boosting & 0.970 & 1.255 & 0.737 & $0.713 \pm 0.018$ & $1.245 \pm 0.034$ \\
        Baseline (CONUS) & Random Forest & 1.287 & 1.631 & 0.556 & $0.520 \pm 0.028$ & $1.609 \pm 0.027$ \\
        Baseline (CONUS) & Gradient Boosting & 1.316 & 1.649 & 0.546 & $0.498 \pm 0.037$ & $1.646 \pm 0.046$ \\
        \bottomrule
    \end{tabularx}
\end{table}

Figure~\ref{fig:environmental_performance_gain} summarizes the Random Forest performance gain from county-only predictors to environmental-only and combined predictor sets. The figure clearly shows that environmental predictors produce the largest improvement in predictive performance.

\begin{figure}[htbp]
    \centering
    \safefigure[width=0.85\textwidth]{figures/notebook5/figure_2_random_forest_environmental_performance_gain.pdf}
    \caption{Random Forest performance gain from county-only predictors to environmental-only and combined predictor sets.}
    \label{fig:environmental_performance_gain}
\end{figure}

The improvement was substantial. With the combined predictor set, the Random Forest achieved a test-set $R^2$ of approximately 0.829, RMSE of approximately 1.011, and MAE of approximately 0.779. Under 5-fold cross-validation, the same model achieved mean $R^2 \approx 0.795$ with standard deviation 0.022, and mean RMSE of approximately 1.049 with standard deviation 0.023. This represents a major improvement over the county-only baseline and shows that wildfire EAL is strongly associated with environmental conditions in addition to county-level exposure and resilience variables.

A particularly important result is that even the \emph{environmental-only} model performed strongly, achieving a test-set $R^2$ of approximately 0.793 and cross-validated $R^2 \approx 0.768$. This indicates that environmental variables carry much of the transferable predictive signal. In other words, once hazard-relevant environmental context is introduced, the model becomes far more informative than a purely socioeconomic county baseline.

\section{Feature Importance in the Combined Model}

Feature-importance analysis of the combined Random Forest model showed that the most influential predictors were primarily environmental. Table~\ref{tab:combined_rf_feature_importance_export} lists the feature-importance values for the combined model after environmental augmentation. The strongest contributors included warm-season precipitation, warm-season vapor pressure deficit, shrub fraction, county area, elevation, building value, population, slope, and warm-season wind speed.

\begin{table}[htbp]
    \centering
    \safetablepdf[width=0.8\textwidth]{tables/notebook5/table_9_combined_random_forest_feature_importance.pdf}
    \caption{Feature importance values for the combined Random Forest model after environmental augmentation.}
    \label{tab:combined_rf_feature_importance_export}
\end{table}

Figure~\ref{fig:combined_rf_feature_importance} visualizes the top feature importances for the combined Random Forest model. The dominance of environmental variables in this figure supports the interpretation that the improved model is learning a wildfire-specific environmental risk signal rather than relying only on demographic or economic scale proxies.

\begin{figure}[htbp]
    \centering
    \safefigure[width=0.86\textwidth]{figures/notebook5/figure_4_top15_combined_rf_feature_importance.pdf}
    \caption{Top feature importances for the combined Random Forest model. Environmental predictors dominate the most important variables.}
    \label{fig:combined_rf_feature_importance}
\end{figure}

The prominence of precipitation and vapor pressure deficit is consistent with the broader wildfire-climate literature. Vapor pressure deficit is closely related to atmospheric drying demand and fuel aridity, both of which influence wildfire activity. Similarly, shrub fraction and slope are plausible wildfire-related factors because they influence fuel structure and fire behavior. These results support the interpretation that the combined model is learning a more wildfire-specific risk signal rather than relying mainly on broad demographic or economic proxies.

\section{State-Wise Geographic Validation}

Although random train/test splits provide a useful first benchmark, they may still place geographically similar counties in both training and testing data. This can lead to optimistic performance estimates when the data are spatially structured. To assess whether the model generalized more robustly across geography, a stricter validation experiment was conducted using grouped cross-validation by state.

Table~\ref{tab:statewise_validation_results_export} summarizes the random split and state-wise grouped validation results for the county-only, environmental-only, and combined feature sets. The grouped validation results are lower than the random-split results, as expected, because the grouped setup requires the model to generalize to held-out states rather than only unseen counties from the same geographic mixture.

\begin{table}[htbp]
    \centering
    \caption{Random split and state-wise grouped validation results for county-only, environmental-only, and combined feature sets.}
    \label{tab:statewise_validation_results_export}
    \scriptsize
    \setlength{\tabcolsep}{4pt}
    \begin{tabularx}{\textwidth}{>{\raggedright\arraybackslash}p{3.4cm} >{\centering\arraybackslash}p{0.9cm} *{6}{>{\centering\arraybackslash}p{1.15cm}}}
        \toprule
        Feature set & $n$ features & Random $R^2$ & Random RMSE & Random MAE & Grouped CV $R^2$ & Grouped CV RMSE & Grouped CV MAE \\
        \midrule
        Combined baseline + environment & 20 & 0.806 & 1.069 & 0.809 & $0.603 \pm 0.041$ & $1.404 \pm 0.155$ & $1.080 \pm 0.113$ \\
        Environmental only & 10 & 0.782 & 1.133 & 0.830 & $0.522 \pm 0.074$ & $1.534 \pm 0.164$ & $1.169 \pm 0.117$ \\
        Baseline (CONUS) & 10 & 0.546 & 1.636 & 1.284 & $0.290 \pm 0.101$ & $1.863 \pm 0.109$ & $1.486 \pm 0.098$ \\
        \bottomrule
    \end{tabularx}
\end{table}

Figure~\ref{fig:random_vs_statewise_r2} compares random-split $R^2$ with state-wise grouped-validation $R^2$. The figure shows that all feature sets experience a performance reduction under grouped validation, but the combined model remains the strongest setup.

\begin{figure}[htbp]
    \centering
    \safefigure[width=0.85\textwidth]{figures/notebook6/figure_1_random_vs_statewise_grouped_r2.pdf}
    \caption{Comparison between random-split $R^2$ and state-wise grouped-validation $R^2$.}
    \label{fig:random_vs_statewise_r2}
\end{figure}

Under this stronger state-wise validation, the combined Random Forest model remained the best-performing setup. Its random-split $R^2$ was approximately 0.806, while its grouped-by-state cross-validated $R^2$ was approximately 0.602. Although this is a noticeable reduction, it is expected because state-wise validation is substantially harder than random splitting. More importantly, the model does not collapse under this stricter evaluation, which indicates that it is learning a meaningful transferable wildfire-risk signal rather than only benefiting from random geographic mixing.

The comparison between feature sets under grouped validation is especially informative. The baseline county-only model achieved grouped $R^2 \approx 0.290$, the environmental-only model achieved grouped $R^2 \approx 0.522$, and the combined model achieved grouped $R^2 \approx 0.602$. This clearly shows that environmental variables provide the strongest transferable predictive signal across states, while county-level exposure and resilience variables still add useful incremental information when combined with them.

Figure~\ref{fig:grouped_model_feature_importance} presents feature importances for the best model under state-wise grouped validation. The important predictors remain consistent with the environmental-augmentation results, reinforcing the importance of climate, land-cover, terrain, and exposure-related information.

\begin{figure}[htbp]
    \centering
    \safefigure[width=0.84\textwidth]{figures/notebook6/figure_2_best_grouped_model_feature_importance.pdf}
    \caption{Feature importances for the best model under state-wise grouped validation.}
    \label{fig:grouped_model_feature_importance}
\end{figure}

Figure~\ref{fig:grouped_predicted_vs_actual} shows predicted versus actual log-transformed wildfire EAL under grouped validation. This diagnostic figure helps illustrate the additional difficulty of predicting held-out states, while still showing that the model retains meaningful predictive structure.

\begin{figure}[htbp]
    \centering
    \safefigure[width=0.65\textwidth]{figures/notebook6/figure_3_grouped_validation_predicted_vs_actual.pdf}
    \caption{Predicted versus actual log-transformed wildfire EAL under grouped validation.}
    \label{fig:grouped_predicted_vs_actual}
\end{figure}

\section{County-Level Risk Assessment Output}
\label{sec:pipeline_output}

The complete pipeline produces a county-level wildfire risk-assessment output rather than only a model-performance score. After the input data are cleaned and merged, each county is represented by a feature vector containing exposure, value, vulnerability, resilience, topographic, land-cover, and climate predictors. The trained model then predicts log-transformed wildfire EAL for each county. These predictions can be converted back into monetary units using the inverse logarithmic transformation:

\begin{equation}
\widehat{\mathrm{wildfire\_eal}} = \exp(\hat{y}) - 1.
\end{equation}

After model evaluation, the final Random Forest pipeline was refit on the complete cleaned combined dataset in order to generate county-level wildfire risk-assessment outputs. This final step is not a separate validation experiment; rather, it represents the operational use of the trained pipeline after the model-selection and validation stages. The resulting predicted EAL values can be used to rank counties, identify high-risk areas, assign percentile-based risk categories, and compare observed and predicted wildfire losses.

Table~\ref{tab:top25_predicted_wildfire_risk_counties} presents the final county-level risk-assessment output from the complete pipeline. The table ranks counties according to predicted wildfire Expected Annual Loss and assigns percentile-based risk categories. This demonstrates that the proposed workflow produces an interpretable wildfire-risk assessment output in monetary units, rather than only reporting model-performance scores.

% THIS TABLE GENRATED USING LATEX BELOW
%\begin{table}[htbp]
    %\centering
    %\safetablepdf[width=\textwidth]{tables/notebook6/table_7_top25_predicted_wildfire_risk_counties.pdf}
    %\caption{Top 25 counties by predicted wildfire Expected Annual Loss from the final county-level risk-assessment pipeline.}
    %\label{tab:top25_predicted_wildfire_risk_counties}
%\end{table}

\begin{table}[htbp]
\centering
\caption{Top 25 counties by predicted wildfire Expected Annual Loss from the final county-level risk-assessment pipeline.}
\label{tab:top25_predicted_wildfire_risk_counties}
\scriptsize
\setlength{\tabcolsep}{3pt}
\renewcommand{\arraystretch}{1.08}
\resizebox{\textwidth}{!}{%
\begin{tabular}{lllrrrrl}
\toprule
State & County & County FIPS & Observed wildfire EAL & Predicted wildfire EAL & Predicted risk percentile & Predicted risk category \\
\midrule
California & Riverside & \texttt{06065} & 346,642,805 & 146,874,987 & 100.0000 & Very High \\
California & San Diego & \texttt{06073} & 430,155,987 & 127,406,725 & 99.9683 & Very High \\
California & Los Angeles & \texttt{06037} & 155,191,812 & 123,291,040 & 99.9366 & Very High \\
California & San Bernardino & \texttt{06071} & 145,737,649 & 75,324,177 & 99.9049 & Very High \\
California & Orange & \texttt{06059} & 57,585,624 & 73,917,750 & 99.8732 & Very High \\
Nevada & Clark & \texttt{32003} & 24,456,604 & 41,937,132 & 99.8415 & Very High \\
Arizona & Maricopa & \texttt{04013} & 42,188,453 & 38,969,217 & 99.8098 & Very High \\
California & Ventura & \texttt{06111} & 50,107,804 & 38,907,335 & 99.7781 & Very High \\
California & Kern & \texttt{06029} & 38,483,037 & 36,200,689 & 99.7464 & Very High \\
California & El Dorado & \texttt{06017} & 55,205,076 & 34,812,071 & 99.7146 & Very High \\
California & Shasta & \texttt{06089} & 35,386,458 & 27,840,399 & 99.6829 & Very High \\
Colorado & Jefferson & \texttt{08059} & 48,732,848 & 27,386,437 & 99.6512 & Very High \\
Utah & Washington & \texttt{49053} & 67,561,249 & 26,959,787 & 99.6195 & Very High \\
Nevada & Washoe & \texttt{32031} & 26,514,858 & 24,483,974 & 99.5878 & Very High \\
Utah & Utah & \texttt{49049} & 35,012,802 & 24,338,860 & 99.5561 & Very High \\
California & Placer & \texttt{06061} & 25,175,654 & 24,061,520 & 99.5244 & Very High \\
Arizona & Pima & \texttt{04019} & 37,561,739 & 23,622,923 & 99.4927 & Very High \\
Oregon & Jackson & \texttt{41029} & 25,011,497 & 22,517,245 & 99.4610 & Very High \\
Utah & Salt Lake & \texttt{49035} & 25,182,533 & 21,831,022 & 99.4293 & Very High \\
California & Calaveras & \texttt{06009} & 30,138,704 & 21,459,704 & 99.3976 & Very High \\
Colorado & Douglas & \texttt{08035} & 25,736,383 & 20,727,496 & 99.3659 & Very High \\
Arizona & Yavapai & \texttt{04025} & 28,707,460 & 20,548,261 & 99.3342 & Very High \\
Arizona & Coconino & \texttt{04005} & 28,503,633 & 19,925,838 & 99.3025 & Very High \\
California & Santa Barbara & \texttt{06083} & 23,974,974 & 19,512,610 & 99.2708 & Very High \\
California & Tuolumne & \texttt{06109} & 27,257,984 & 19,385,286 & 99.2391 & Very High \\
\bottomrule
\end{tabular}%
}
\end{table}

The table should be interpreted as the final assessment product of the FEMA/NRI-based modeling workflow. Counties with the highest predicted wildfire EAL represent locations where the combined exposure, vulnerability, resilience, terrain, land-cover, and climate predictors indicate the greatest expected wildfire-loss burden. The percentile-based risk categories provide a simple way to translate continuous monetary predictions into interpretable risk classes.

This output is important for the thesis because it shows that the proposed approach is not limited to wildfire occurrence prediction. Instead, it produces a consequence-oriented risk estimate in monetary units. The pipeline therefore connects the literature gap identified in Chapter~\ref{chap:litreview} with the experimental results presented in this chapter: wildfire hazard-related environmental information is combined with exposure and vulnerability indicators to estimate expected wildfire losses.

\section{Exploratory Ecological Risk Proxy}
\label{sec:ecological_proxy}

The main FEMA/NRI model predicts monetary wildfire Expected Annual Loss, but the thesis title and literature review also emphasize ecological wildfire consequences. To connect the empirical pipeline with this ecological dimension, an exploratory ecological risk proxy was constructed using the environmental predictors already available in the county-level dataset. This analysis does not estimate ecological loss directly. Instead, it identifies counties where high predicted wildfire risk overlaps with high natural vegetation exposure.

Table~\ref{tab:ecological_proxy_definition} summarizes the proxy construction logic. Natural vegetation exposure is represented by the combined share of forest, shrub, and grass land cover. This value is then multiplied by the predicted wildfire-risk percentile from the final FEMA/NRI model to create an ecological risk proxy. Higher values indicate counties where predicted wildfire risk and natural vegetation exposure are both high.

%This was generated using LAtex code
%\begin{table}[htbp]
    %\centering
    %\safetablepdf[width=0.9\textwidth]{tables/notebook7/table_1_ecological_proxy_definition.pdf}
    %\caption{Definition of the exploratory ecological risk proxy.}
    %\label{tab:ecological_proxy_definition}
%\end{table}

\begin{table}[htbp]
\centering
\caption{Definition of the exploratory ecological risk proxy.}
\label{tab:ecological_proxy_definition}
\scriptsize
\setlength{\tabcolsep}{4pt}
\renewcommand{\arraystretch}{1.15}
\begin{tabularx}{\textwidth}{L{3.4cm} L{5.7cm} X}
\toprule
Component & Definition & Interpretation \\
\midrule
Forest fraction & County fraction covered by forest land cover & Natural vegetation exposure \\
Shrub fraction & County fraction covered by shrub land cover & Natural vegetation exposure \\
Grass fraction & County fraction covered by grass land cover & Natural vegetation exposure \\
Natural vegetation fraction & Forest fraction + shrub fraction + grass fraction & Approximate natural ecosystem exposure at county level \\
Predicted risk percentile & Percentile rank of predicted wildfire EAL & Economic wildfire risk signal from the final FEMA/NRI model \\
Ecological risk proxy & Predicted risk percentile multiplied by natural vegetation fraction & Overlap of predicted wildfire risk and natural vegetation exposure \\
Ecological risk category & Percentile-based category of ecological risk proxy & Screening class for exploratory ecological wildfire-risk assessment \\
\bottomrule
\end{tabularx}
\end{table}

Figure~\ref{fig:predicted_eal_vs_natural_vegetation} shows the relationship between predicted wildfire EAL and natural vegetation fraction. The figure helps distinguish counties with high monetary wildfire risk from counties where wildfire risk overlaps with larger natural vegetation exposure. Figure~\ref{fig:ecological_risk_proxy_distribution} shows the distribution of the ecological risk proxy across counties.

\begin{figure}[htbp]
    \centering
    \safefigure[width=0.82\textwidth]{figures/notebook7/figure_1_predicted_eal_vs_natural_vegetation.pdf}
    \caption{Relationship between predicted wildfire Expected Annual Loss and natural vegetation fraction.}
    \label{fig:predicted_eal_vs_natural_vegetation}
\end{figure}

\begin{figure}[htbp]
    \centering
    \safefigure[width=0.82\textwidth]{figures/notebook7/figure_2_ecological_risk_proxy_distribution.pdf}
    \caption{Distribution of the exploratory ecological risk proxy across counties.}
    \label{fig:ecological_risk_proxy_distribution}
\end{figure}

Table~\ref{tab:top25_ecological_risk_proxy_counties} presents the counties with the highest ecological risk proxy values. These counties should not be interpreted as having measured ecological losses. Rather, they represent locations where the model predicts high wildfire-loss risk and where natural vegetation exposure is also high, making them useful candidates for further ecological vulnerability or recovery analysis. The natural vegetation fraction is computed as the sum of forest, shrub, and grass land-cover fractions, as defined in Table~\ref{tab:ecological_proxy_definition}; supporting component values for the top ecological-risk counties are provided in Appendix Table~\ref{tab:app_n7_ecological_component_values}.

% This table generated through latex
%\begin{table}[htbp]
    %\centering
    %\safetablepdf[width=\textwidth]{tables/notebook7/table_2_top25_ecological_risk_proxy_counties.pdf}
    %\caption{Top 25 counties by exploratory ecological risk proxy.}
    %\label{tab:top25_ecological_risk_proxy_counties}
%\end{table}

\begin{longtable}{p{1.8cm} p{1.9cm} c r r r r}
\caption{Top 25 counties by exploratory ecological risk proxy. All listed counties fall in the Very High ecological risk category.}
\label{tab:top25_ecological_risk_proxy_counties}\\

\toprule
State & County & FIPS &
\shortstack{Natural\\vegetation\\fraction} &
\shortstack{Predicted\\wildfire EAL} &
\shortstack{Predicted\\risk\\percentile} &
\shortstack{Ecological\\risk proxy} \\
\midrule
\endfirsthead

\multicolumn{7}{c}{\tablename~\thetable{} -- continued from previous page} \\
\toprule
State & County & FIPS &
\shortstack{Natural\\vegetation\\fraction} &
\shortstack{Predicted\\wildfire EAL} &
\shortstack{Predicted\\risk\\percentile} &
\shortstack{Ecological\\risk proxy} \\
\midrule
\endhead

\midrule
\multicolumn{7}{r}{Continued on next page} \\
\endfoot

\bottomrule
\endlastfoot

New Mexico & Lincoln & \texttt{35027} & 0.9904 & 7,664,865 & 97.98 & 97.04 \\
Arizona & Yavapai & \texttt{04025} & 0.9765 & 20,276,267 & 99.37 & 97.04 \\
Nevada & Elko & \texttt{32007} & 0.9715 & 15,797,063 & 99.05 & 96.23 \\
Arizona & Coconino & \texttt{04005} & 0.9690 & 20,168,753 & 99.30 & 96.22 \\
Arizona & Navajo & \texttt{04017} & 0.9762 & 9,334,416 & 98.29 & 95.96 \\
New Mexico & Grant & \texttt{35017} & 0.9820 & 5,393,697 & 97.47 & 95.71 \\
Arizona & Gila & \texttt{04007} & 0.9774 & 7,569,287 & 97.91 & 95.70 \\
New Mexico & San Miguel & \texttt{35047} & 0.9853 & 3,982,625 & 96.93 & 95.51 \\
Utah & Washington & \texttt{49053} & 0.9556 & 28,360,770 & 99.68 & 95.26 \\
Idaho & Boise & \texttt{16015} & 0.9837 & 3,620,263 & 96.68 & 95.11 \\
California & Mariposa & \texttt{06043} & 0.9712 & 6,893,095 & 97.82 & 95.00 \\
Arizona & Apache & \texttt{04001} & 0.9792 & 3,428,775 & 96.59 & 94.58 \\
Arizona & Mohave & \texttt{04015} & 0.9584 & 11,952,896 & 98.58 & 94.48 \\
Arizona & Santa Cruz & \texttt{04023} & 0.9672 & 5,884,929 & 97.63 & 94.43 \\
Idaho & Idaho & \texttt{16049} & 0.9602 & 8,220,216 & 98.17 & 94.26 \\
New Mexico & Otero & \texttt{35035} & 0.9568 & 11,682,361 & 98.51 & 94.25 \\
New Mexico & Catron & \texttt{35003} & 0.9943 & 1,942,928 & 94.63 & 94.08 \\
New Mexico & Colfax & \texttt{35007} & 0.9711 & 3,727,119 & 96.81 & 94.01 \\
Montana & Ravalli & \texttt{30081} & 0.9578 & 7,982,301 & 98.07 & 93.94 \\
New Mexico & Sandoval & \texttt{35043} & 0.9645 & 4,516,141 & 97.19 & 93.74 \\
Idaho & Adams & \texttt{16003} & 0.9789 & 2,632,951 & 95.54 & 93.53 \\
California & Calaveras & \texttt{06009} & 0.9400 & 21,267,869 & 99.43 & 93.47 \\
Oregon & Jackson & \texttt{41029} & 0.9394 & 21,874,514 & 99.46 & 93.44 \\
Utah & Iron & \texttt{49021} & 0.9503 & 9,140,679 & 98.26 & 93.38 \\
Arizona & Pima & \texttt{04019} & 0.9382 & 23,527,659 & 99.49 & 93.35 \\

\end{longtable}



Figure~\ref{fig:economic_ecological_risk_quadrants_notebook7} provides an additional visualization of the relationship between predicted economic wildfire risk and natural vegetation exposure. This supports the interpretation that monetary risk and ecological exposure are related but not identical dimensions of wildfire consequence assessment.

\begin{figure}[htbp]
    \centering
    \safefigure[width=0.82\textwidth]{figures/notebook7/figure_3_economic_vs_ecological_risk_quadrants.pdf}
    \caption{Economic wildfire-risk percentile and natural vegetation exposure used to motivate ecological risk screening.}
    \label{fig:economic_ecological_risk_quadrants_notebook7}
\end{figure}

\section{Economic-Ecological Risk Overlap Analysis}
\label{sec:economic_ecological_overlap}

The ecological proxy analysis was extended by grouping counties according to the overlap between predicted monetary wildfire risk and natural vegetation exposure. This step provides a simple way to separate counties where both dimensions are high from counties where only monetary risk or ecological exposure is high. The analysis is exploratory, but it helps clarify how the FEMA/NRI monetary-risk model can be connected to ecological screening without claiming to directly estimate ecological damage.

Table~\ref{tab:economic_ecological_quadrant_summary} summarizes the four economic--ecological overlap groups. Counties in the high economic risk and high ecological exposure group are especially important because they combine high predicted wildfire-loss risk with high natural vegetation exposure. Counties with lower economic risk but high ecological exposure may also deserve ecological monitoring, even if their predicted monetary losses are lower.

% This also generated through latex
%\begin{table}[htbp]
    %\centering
    %\safetablepdf[width=\textwidth]{tables/notebook8/table_1_economic_ecological_quadrant_summary.pdf}
    %\caption{Economic--ecological risk-overlap quadrant summary.}
    %\label{tab:economic_ecological_quadrant_summary}
%\end{table}

\begin{table}[htbp]
\centering
\caption{Economic--ecological risk-overlap quadrant summary.}
\label{tab:economic_ecological_quadrant_summary}
\scriptsize
\setlength{\tabcolsep}{3pt}
\renewcommand{\arraystretch}{1.15}
\resizebox{\textwidth}{!}{%
\begin{tabular}{lrrrrrr}
\toprule
Risk-overlap quadrant & County count & Mean observed wildfire EAL & Mean predicted wildfire EAL & Mean predicted risk percentile & Mean natural vegetation fraction & Mean ecological risk proxy \\
\midrule
High economic risk + high ecological exposure & 238 & 5,315,258 & 3,466,650 & 91.37 & 0.94 & 85.62 \\
High economic risk + lower ecological exposure & 395 & 6,047,578 & 3,719,458 & 89.19 & 0.55 & 49.63 \\
Lower economic risk + high ecological exposure & 395 & 88,208 & 66,091 & 44.94 & 0.92 & 41.48 \\
Lower economic risk + lower ecological exposure & 2,135 & 59,641 & 49,950 & 39.10 & 0.46 & 19.47 \\
\bottomrule
\end{tabular}%
}
\end{table}

Figure~\ref{fig:economic_ecological_quadrant_scatter} visualizes the same overlap logic. The scatter plot separates counties by predicted wildfire-risk percentile and natural vegetation fraction, highlighting the counties where economic risk and ecological exposure coincide.

\begin{figure}[htbp]
    \centering
    \safefigure[width=0.82\textwidth]{figures/notebook8/figure_1_economic_ecological_quadrant_scatter.pdf}
    \caption{Economic--ecological overlap between predicted wildfire-risk percentile and natural vegetation exposure.}
    \label{fig:economic_ecological_quadrant_scatter}
\end{figure}

Table~\ref{tab:high_economic_high_ecological_counties} lists the top counties in the high economic risk and high ecological exposure group. This table provides a practical screening output for identifying counties where further ecological vulnerability analysis, burn-severity assessment, biodiversity exposure mapping, or recovery monitoring may be especially relevant.

%This was also generated with LATEX
%\begin{table}[htbp]
    %\centering
    %\safetablepdf[width=\textwidth]{tables/notebook8/table_2_high_economic_high_ecological_counties.pdf}
    %\caption{Top counties with high predicted economic wildfire risk and high natural vegetation exposure.}
    %\label{tab:high_economic_high_ecological_counties}
%\end{table}

\begin{table}[htbp]
\centering
\caption{Top counties with high predicted economic wildfire risk and high natural vegetation exposure.}
\label{tab:high_economic_high_ecological_counties}
\tiny
\setlength{\tabcolsep}{2.8pt}
\renewcommand{\arraystretch}{1.06}
\resizebox{\textwidth}{!}{%
\begin{tabular}{lllrrrrr}
\toprule
State & County & County FIPS & Predicted wildfire EAL & Predicted risk percentile & Natural vegetation fraction & Ecological risk proxy & Ecological risk percentile \\
\midrule
New Mexico & Lincoln & \texttt{35027} & 7,664,865 & 97.98 & 0.99 & 97.04 & 100.00 \\
Arizona & Yavapai & \texttt{04025} & 20,276,267 & 99.37 & 0.98 & 97.04 & 99.97 \\
Nevada & Elko & \texttt{32007} & 15,797,063 & 99.05 & 0.97 & 96.23 & 99.94 \\
Arizona & Coconino & \texttt{04005} & 20,168,753 & 99.30 & 0.97 & 96.22 & 99.91 \\
Arizona & Navajo & \texttt{04017} & 9,334,416 & 98.29 & 0.98 & 95.96 & 99.87 \\
New Mexico & Grant & \texttt{35017} & 5,393,697 & 97.47 & 0.98 & 95.71 & 99.84 \\
Arizona & Gila & \texttt{04007} & 7,569,287 & 97.91 & 0.98 & 95.70 & 99.81 \\
New Mexico & San Miguel & \texttt{35047} & 3,982,625 & 96.93 & 0.99 & 95.51 & 99.78 \\
Utah & Washington & \texttt{49053} & 28,360,770 & 99.68 & 0.96 & 95.26 & 99.75 \\
Idaho & Boise & \texttt{16015} & 3,620,263 & 96.68 & 0.98 & 95.11 & 99.72 \\
California & Mariposa & \texttt{06043} & 6,893,095 & 97.82 & 0.97 & 95.00 & 99.68 \\
Arizona & Apache & \texttt{04001} & 3,428,775 & 96.59 & 0.98 & 94.58 & 99.65 \\
Arizona & Mohave & \texttt{04015} & 11,952,896 & 98.58 & 0.96 & 94.48 & 99.62 \\
Arizona & Santa Cruz & \texttt{04023} & 5,884,929 & 97.63 & 0.97 & 94.43 & 99.59 \\
Idaho & Idaho & \texttt{16049} & 8,220,216 & 98.17 & 0.96 & 94.26 & 99.56 \\
New Mexico & Otero & \texttt{35035} & 11,682,361 & 98.51 & 0.96 & 94.25 & 99.53 \\
New Mexico & Catron & \texttt{35003} & 1,942,928 & 94.63 & 0.99 & 94.08 & 99.49 \\
New Mexico & Colfax & \texttt{35007} & 3,727,119 & 96.81 & 0.97 & 94.01 & 99.46 \\
Montana & Ravalli & \texttt{30081} & 7,982,301 & 98.07 & 0.96 & 93.94 & 99.43 \\
New Mexico & Sandoval & \texttt{35043} & 4,516,141 & 97.19 & 0.96 & 93.74 & 99.40 \\
Idaho & Adams & \texttt{16003} & 2,632,951 & 95.54 & 0.98 & 93.53 & 99.37 \\
California & Calaveras & \texttt{06009} & 21,267,869 & 99.43 & 0.94 & 93.47 & 99.34 \\
Oregon & Jackson & \texttt{41029} & 21,874,514 & 99.46 & 0.94 & 93.44 & 99.30 \\
Utah & Iron & \texttt{49021} & 9,140,679 & 98.26 & 0.95 & 93.38 & 99.27 \\
Arizona & Pima & \texttt{04019} & 23,527,659 & 99.49 & 0.94 & 93.35 & 99.24 \\
\bottomrule
\end{tabular}%
}
\end{table}

\section{Supporting Results from the Vietnam Wildfire Prediction Study}
\label{sec:vietnam_results}

Although the main empirical pipeline of this thesis focuses on FEMA/NRI county-level wildfire Expected Annual Loss modeling in the United States, a previous supporting study was also conducted for An Giang Province, Vietnam. This study is included here to demonstrate earlier fine-scale wildfire prediction work and to connect the present loss-assessment pipeline with the author's prior experience in deep-learning-based wildfire risk modeling.

The Vietnam study developed an enhanced deep neural network for classifying significant wildfire events in An Giang Province. The model used multi-source environmental inputs, including vegetation indices, land-use/land-cover information, MERRA-2 weather variables, and dynamic fire-growth metrics. The binary target was based on significant fire events, with a burned-area threshold of 0.6~km$^2$. The model also used class weighting, threshold optimization, and SHAP-based interpretation to improve operational usefulness and explainability \citep{suliman2026angiang}.

Table~\ref{tab:vietnam_study_results} summarizes the key performance results from the Vietnam study. The enhanced DNN achieved 91.4\% overall accuracy, 74.8\% recall for significant fires, 92.5\% precision, 82.7\% F1-score, 97.7\% specificity, and a 2.3\% false-positive rate. These results indicate that the model achieved a useful balance between detecting significant fires and limiting false alarms.

\begin{table}[htbp]
\centering
\caption{Key results from the supporting An Giang, Vietnam wildfire prediction study.}
\label{tab:vietnam_study_results}
\begin{tabular}{lc}
\toprule
Metric & Value \\
\midrule
Overall accuracy & 91.4\% \\
Significant fire recall & 74.8\% \\
Significant fire precision & 92.5\% \\
F1-score & 82.7\% \\
Specificity & 97.7\% \\
False positive rate & 2.3\% \\
False negative rate & 25.2\% \\
Dominant SHAP feature & Fire spread rate \\
Fire spread rate importance & 37.6\% \\
Optimal decision threshold & 0.65 \\
\bottomrule
\end{tabular}
\end{table}

The Vietnam results are not treated as part of the FEMA/NRI county-level loss model because they address a different task: significant-fire classification rather than monetary Expected Annual Loss prediction. Nevertheless, they are relevant to the thesis because they show how wildfire risk modeling can be improved using dynamic fire-growth features, multi-source environmental data, threshold-sensitive evaluation, and model interpretability. In this sense, the Vietnam study provides a fine-scale prediction-oriented complement to the main FEMA/NRI loss-assessment pipeline.

\section{Summary of Current Results}

Table~\ref{tab:prelim_results_summary} summarizes the current main experiments. The table condenses the main result of the chapter: the county-only baseline provides a meaningful starting point, environmental variables substantially improve predictive performance, and the combined model remains the strongest setup even under state-wise validation.

\begin{table}[htbp]
\centering
\caption{Summary of county-level wildfire EAL modeling results.}
\label{tab:prelim_results_summary}
\resizebox{\textwidth}{!}{%
\begin{tabular}{lcccc}
\hline
Model setup & Best model & Test $R^2$ & CV / grouped $R^2$ & Notes \\
\hline
County-only baseline & Random Forest & 0.483 & 0.504 & Accessible county variables only \\
Environmental-only (CONUS) & Random Forest & 0.793 & 0.768 & Topography, land cover, climate \\
Combined baseline + environment (CONUS) & Random Forest & 0.829 & 0.795 & Standard random cross-validation \\
Combined baseline + environment (state-wise) & Random Forest & 0.806 & 0.602 & Grouped cross-validation by state \\
\hline
\end{tabular}%
}
\end{table}

Overall, these results support three main conclusions. First, county-level wildfire expected annual loss can be modeled with meaningful predictive power using accessible public datasets. Second, environmental variables contribute much more strongly than county socioeconomic variables alone, especially under geographically stricter validation. Third, the combined model remains reasonably robust even when evaluated across held-out states, which makes it a credible current baseline for the thesis.

\section{Interpretation and Implications}

The experimental results demonstrate that moving from hazard-oriented wildfire modeling to loss-oriented wildfire-risk assessment is feasible using publicly accessible datasets. The FEMA/NRI-based pipeline shows that county-level wildfire Expected Annual Loss can be modeled with meaningful predictive power when exposure, value, vulnerability, resilience, topographic, land-cover, and climate variables are combined in a structured machine-learning workflow. This is important for the thesis because it provides a consequence-oriented wildfire-risk modeling framework in monetary units rather than a model limited only to wildfire occurrence or burned area.

The results also show that environmental predictors substantially improve model performance compared with county-level socioeconomic and exposure variables alone. In the county-only baseline, Random Forest achieved moderate predictive performance, indicating that exposure, value, vulnerability, and resilience variables contain useful but incomplete information. After environmental augmentation, performance increased substantially, and the combined model remained the strongest setup under both random validation and state-wise grouped validation. This supports the central argument of the thesis: wildfire losses are shaped by both the exposed human/economic system and the environmental conditions that influence wildfire hazard.

The state-wise grouped validation results are especially important because they provide a more demanding test of geographic generalization. Although performance declined compared with random splitting, the combined model did not collapse when evaluated across held-out states. This suggests that the pipeline captures a transferable wildfire-risk signal rather than only benefiting from random geographic mixing. At the same time, the reduction in grouped-validation performance shows that regional generalization remains challenging and should be considered when interpreting the model as a decision-support tool.

The final county-level risk-assessment output further strengthens the practical value of the pipeline. By converting predicted log-transformed wildfire EAL back into monetary units and ranking counties by predicted wildfire loss, the model produces an interpretable risk-assessment product. The exploratory ecological proxy and economic--ecological overlap analysis add a screening layer that identifies where high predicted wildfire risk coincides with high natural vegetation exposure. This makes the proposed workflow more useful than a purely technical regression experiment, because the output can support county-level comparison, prioritization, and further investigation of areas with high expected wildfire losses or ecological exposure.

\section{Limitations and Future Work}

Several limitations should be acknowledged. First, the target variable is county-level wildfire Expected Annual Loss from the FEMA National Risk Index. Although this provides a useful and publicly accessible monetary risk target, it is spatially aggregated and does not capture event-level losses, parcel-level exposure, or the detailed mechanisms through which individual wildfire events cause damage. Therefore, the model should be interpreted as a county-level risk-assessment framework rather than an event-level damage prediction system.

Second, the environmental augmentation is restricted to the contiguous United States because the selected land-cover and climate variables were prepared within a CONUS-compatible workflow. This provides a consistent national-scale modeling setup, but it also means that the current environmental feature set does not directly cover non-CONUS areas or all county-equivalent units. A small number of unmatched county-equivalent units were excluded during the environmental merge.

Third, although state-wise grouped validation provides a stronger test than random splitting, geographic transferability remains imperfect. The reduction in performance under grouped validation indicates that wildfire-loss relationships vary across regions due to differences in climate, vegetation, settlement patterns, suppression systems, exposure, and local fire regimes. Future work should therefore explore additional spatial validation designs, regional calibration strategies, and area-of-applicability analysis.

Fourth, the current model uses Random Forest feature importance as an initial interpretation tool. While this is useful for identifying influential predictors, future work could incorporate more detailed explainability methods such as SHAP values to better understand nonlinear interactions among exposure, vulnerability, land cover, topography, and climate variables.

Fifth, the empirical model does not directly predict ecological losses as a separate target variable. The exploratory ecological proxy and economic--ecological overlap analysis provide a screening perspective based on natural vegetation exposure and predicted wildfire risk, but they should not be interpreted as measured ecological damage, burn severity, biodiversity loss, or ecosystem-service loss. Future work should therefore connect this screening layer with explicit ecological vulnerability, recovery, or ecosystem-value datasets.

Future extensions can proceed in several directions. The first direction is methodological refinement, including additional wildfire-relevant predictors, improved spatial validation, uncertainty quantification, and county-level error analysis. The second direction is data enrichment, especially through more detailed event-level loss data, building-level exposure layers, infrastructure data, smoke-health impacts, suppression-cost data, and ecological vulnerability indicators. The third direction is fine-scale modeling, where the exploratory Creek Fire tensor-based workflow described in Appendix~\ref{app:creek_fire_pipeline} could be used as a future extension toward spatially detailed wildfire-risk modeling. Together, these extensions would help move from the current county-level loss-assessment framework toward a more comprehensive multi-scale wildfire-risk and wildfire-loss modeling system.
\chapter{Discussion}
\label{chap:discussion}

The experimental results demonstrate that county-level wildfire Expected Annual Loss can be modeled with meaningful predictive power using accessible public data. A key finding is that environmental augmentation substantially improves performance over the county-only baseline. This suggests that loss-oriented wildfire modeling should not rely only on exposure, value, vulnerability, and socioeconomic indicators, but should also include hazard-relevant environmental conditions such as precipitation, vapor pressure deficit, vegetation composition, slope, elevation, and wind.

Unlike previous wildfire-risk studies that often evaluate predictive accuracy without translating results into consequence-oriented outputs, the present thesis converts model predictions into county-level wildfire Expected Annual Loss estimates, risk rankings, and exploratory ecological screening outputs.

The results support the central pipeline logic of the thesis. The FEMA National Risk Index provides a monetary wildfire-loss target, while the county-level predictors represent exposure, asset value, vulnerability, and resilience. Environmental augmentation then adds wildfire-hazard context derived from topography, land cover, and climate. Finally, machine-learning regression models and geographic validation convert these inputs into a county-level wildfire risk prediction and assessment workflow. In this way, the thesis moves beyond wildfire occurrence prediction and produces a consequence-oriented risk estimate in monetary units.

A central result is the improvement from the county-only baseline to the combined baseline-plus-environment model. The county-only model shows that exposure, value, vulnerability, and resilience variables contain useful predictive information. However, the much stronger performance of the environmental-only and combined models indicates that wildfire EAL is strongly shaped by environmental conditions. This supports the broader argument that wildfire losses should be understood as the result of interactions among hazard, exposure, vulnerability, and resilience rather than as a function of exposed assets alone.

The state-wise grouped validation provides an important perspective on geographic generalization. Although model performance decreases when entire states are held out, the combined model still retains substantial predictive power. This indicates that the model is learning a meaningful transferable wildfire-risk signal rather than merely benefiting from random geographic similarity between training and test observations. At the same time, the reduction in grouped-validation performance shows that regional generalization remains challenging. County-level wildfire-loss relationships vary across regions because of differences in fire regimes, vegetation structure, climate, settlement patterns, suppression systems, and local socioeconomic context.

The final risk-assessment output further strengthens the practical interpretation of the pipeline. By converting predicted log-transformed wildfire EAL back into monetary units, the model can rank counties by predicted wildfire loss and assign percentile-based risk categories. This step is important because it shows that the proposed workflow is not only a model-performance exercise. Instead, it produces an interpretable county-level risk-assessment product that can support comparison, prioritization, and further investigation of high-risk areas.

The exploratory ecological risk proxy adds a complementary screening layer to the monetary EAL model. By combining predicted wildfire-risk percentile with natural vegetation exposure, the analysis identifies counties where high predicted wildfire risk overlaps with high forest, shrub, and grass cover. This does not convert the thesis into a direct ecological-loss model, but it strengthens the ecological dimension by showing how the main FEMA/NRI pipeline can be extended toward ecological vulnerability screening. The economic--ecological overlap analysis further distinguishes counties with both high monetary risk and high ecological exposure from counties where only one of these dimensions is high.

Another important implication is conceptual. The experiments support the broader thesis argument that wildfire consequence assessment can be approached as a combined hazard--exposure--vulnerability problem. The county-only baseline captures part of the exposure and vulnerability signal, but the environmental augmentation shows that hazard-relevant climate, terrain, and land-cover information is essential for stronger loss-oriented prediction. This aligns with the literature gap identified in Chapter~\ref{chap:litreview}: wildfire-risk prediction becomes more useful when it is connected to expected consequences rather than treated as an endpoint by itself.

The thesis also contains an earlier Creek Fire tensor-based workflow as supporting work. This workflow was useful for developing a fine-scale data-engineering approach for wildfire modeling, including fire detections, fire perimeters, static geospatial predictors, and time-varying climate variables. However, it is not the main M.S. pipeline because it is not directly connected to monetary loss assessment. Therefore, it is placed in Appendix~\ref{app:creek_fire_pipeline} as a baseline approach that can support future fine-scale extensions of the FEMA/NRI loss-oriented pipeline.

The supporting Vietnam study provides an additional perspective on the broader wildfire-risk modeling direction of this thesis. While the FEMA/NRI pipeline focuses on county-level monetary Expected Annual Loss, the An Giang study focused on significant-fire classification using a compact deep neural network and multi-source environmental predictors. Its results showed strong classification performance, including 91.4\% accuracy, 74.8\% recall, 92.5\% precision, and 97.7\% specificity \citep{suliman2026angiang}. The study also showed that fire spread rate was the dominant SHAP feature, indicating the importance of dynamic fire-growth information for fine-scale wildfire prediction. Together, the Vietnam and FEMA/NRI results show two complementary directions: fine-scale wildfire event prediction and county-level monetary risk/loss assessment.

Despite the strengths of the FEMA/NRI pipeline, several limitations remain. First, the target variable is county-level Expected Annual Loss, which is spatially aggregated and does not represent event-level, parcel-level, or asset-level losses directly. Second, the environmental augmentation is restricted to the contiguous United States and excludes a small number of county-equivalent units because of geometry and dataset compatibility constraints. Third, the current modeling framework focuses on point prediction and feature importance; it does not yet include uncertainty quantification, calibrated prediction intervals, or probabilistic risk categories. Fourth, Random Forest feature importance provides useful first-order interpretation, but more detailed explainability methods such as SHAP could provide a deeper understanding of nonlinear feature interactions.

Although state-wise grouped validation is stricter than random splitting, it is not the only possible spatial validation strategy. Counties in different states can still share similar ecoregions, climate regimes, or vegetation types, while counties within large states may be highly heterogeneous. Future work could therefore evaluate transferability using ecoregion-based, climate-zone-based, or spatial-block validation strategies.

A further limitation concerns the ecological-loss dimension of the thesis. While ecological losses are discussed in the literature review and explored through the proxy-based screening layer, the empirical model does not directly predict ecological damage as a supervised target variable. The implemented pipeline focuses on monetary wildfire Expected Annual Loss from FEMA NRI. Environmental predictors such as forest fraction, shrub fraction, grass fraction, elevation, slope, precipitation, wind speed, and vapor pressure deficit provide ecological and hazard-related context, while the ecological proxy identifies overlap between predicted wildfire risk and natural vegetation exposure. However, this still does not constitute a direct ecological-loss model.

These limitations point toward several future extensions. Future work can improve geographic robustness, extend the environmental predictor set, add uncertainty quantification, and investigate richer representations of wildfire loss using event-level, asset-level, infrastructure, and ecological data. In particular, future ecological-loss modeling should incorporate explicit ecological vulnerability or recovery indicators, such as burn severity, ecosystem value, biodiversity exposure, protected-area overlap, vegetation recovery time, or post-fire resilience metrics. Nevertheless, the present results establish a defendable county-level wildfire loss-assessment pipeline. They show that publicly accessible exposure, vulnerability, resilience, topographic, land-cover, and climate datasets can be combined to produce meaningful wildfire-loss predictions in monetary units and exploratory ecological risk-screening outputs.

\chapter{Conclusion and Future Work}
\label{chap:conclusion}

This thesis investigated the problem of connecting wildfire-risk prediction with wildfire-loss assessment. The central motivation was that much of the recent wildfire literature remains focused on hazard prediction, including ignition likelihood, susceptibility, spread, detection, and fire-danger estimation, while fewer studies translate predictive outputs into economic or ecological consequence estimates. To address this gap, the thesis developed a county-level wildfire-loss modeling framework using publicly accessible datasets.

The literature review showed that recent research has made substantial progress in wildfire prediction, remote sensing, deep learning, ecological vulnerability assessment, and AI-supported decision systems. At the same time, it highlighted that integrated wildfire-loss modeling remains comparatively underdeveloped. This is especially true for scalable frameworks that connect hazard-relevant environmental variables with monetary loss-oriented targets and geographic validation.

The main methodological contribution of the thesis is a complete FEMA/NRI-based county-level wildfire risk prediction and assessment pipeline. The pipeline begins with a county-level wildfire Expected Annual Loss target in monetary units. It then adds county-level exposure, value, vulnerability, and resilience predictors, followed by environmental predictors derived from topography, land cover, and warm-season climate summaries. The resulting feature table is modeled using regression algorithms and evaluated under both random and state-wise grouped validation.

The experimental results show that the county-only baseline provides a meaningful but limited predictive signal. Environmental augmentation produces a substantial improvement, indicating that wildfire losses are strongly associated with terrain, land cover, and climate conditions in addition to county-level exposure and resilience variables. The strongest performance was obtained with the combined Random Forest model. State-wise grouped validation further showed that although geographic generalization is more difficult than random splitting, the combined model remains useful when evaluated across held-out states.

The final pipeline output demonstrates that the model can be used for county-level wildfire risk assessment, not only for regression benchmarking. By converting predicted log-transformed wildfire EAL back into monetary units, the workflow produces predicted wildfire-loss estimates, county rankings, and percentile-based risk categories. The thesis also adds an exploratory ecological risk-screening layer that combines predicted wildfire-risk percentile with natural vegetation exposure. This makes the thesis contribution consequence-oriented: the output is an interpretable wildfire-risk assessment product in monetary units, complemented by an ecological exposure-screening perspective, rather than only a wildfire occurrence or hazard prediction.

The thesis also documented an earlier Creek Fire tensor-based workflow as a supporting baseline and future extension. This workflow demonstrates fine-scale data engineering using fire detections, a fire perimeter, static geospatial predictors, and time-varying climate features. However, the main defendable M.S. contribution is the FEMA/NRI county-level wildfire risk/loss assessment pipeline, because it is directly connected to monetary wildfire-loss estimation.

Several limitations remain. The FEMA NRI wildfire EAL target is spatially aggregated and does not capture event-level or asset-level damage mechanisms directly. The environmental augmentation currently focuses on the contiguous United States and does not fully address all county-equivalent units. The validation strategy includes state-wise grouped validation, but additional spatial validation designs and uncertainty analysis would strengthen the model further. In addition, future work should incorporate more detailed explainability methods, richer loss data, and ecological consequence indicators.

Future research can extend this work in three main directions. First, the modeling framework can be improved with additional wildfire-relevant predictors, uncertainty quantification, regional calibration, and more advanced explainability methods. Second, the data foundation can be enriched using event-level wildfire losses, building-level exposure, infrastructure networks, suppression-cost data, smoke-health impacts, and ecological vulnerability indicators. Third, the exploratory Creek Fire workflow can be developed into a fine-scale spatio-temporal modeling pipeline that complements the county-level FEMA/NRI framework.

The empirical contribution of the thesis remains primarily economic: it models county-level wildfire Expected Annual Loss in monetary units. Ecological losses are addressed through the literature review, through the inclusion of environmental predictors, and through an exploratory ecological proxy based on natural vegetation exposure. However, ecological losses are not directly estimated as a separate damage target. Future work should extend the current framework by integrating ecological vulnerability, burn severity, vegetation recovery, biodiversity exposure, and ecosystem-service loss indicators.

The supporting Vietnam study further demonstrates that fine-scale wildfire prediction can benefit from dynamic fire-growth features, threshold-aware evaluation, and SHAP-based interpretability, while the FEMA/NRI pipeline extends the thesis toward county-level monetary loss assessment.

Part of this work was successfully presented at DCHPC 2026 through the paper entitled ``Mapping Forest Fire Risks Using Deep Learning: A Case Study in An Giang Province, Vietnam.'' This conference presentation supports the broader research direction of the thesis by demonstrating fine-scale wildfire prediction using deep learning, while the main thesis contribution extends the work toward county-level monetary wildfire-loss assessment. The corresponding certificate is included in Appendix~\ref{app:conference_certificate}, and the associated paper is expected to appear as a Web of Science/Scopus paper.

The main novelty of the thesis lies in the general workflow and insights rather than in the specific territory used for implementation. Unlike previous wildfire-prediction studies that often stop at hazard, susceptibility, or burned-area prediction, this thesis connects wildfire-relevant environmental conditions with exposure, vulnerability, resilience, and monetary expected loss. The new methodological contribution is an integrated wildfire-loss assessment pipeline that combines environmental augmentation, machine-learning regression, geographic validation, county-level risk ranking, and exploratory ecological risk screening.

The key insights gained are threefold. First, environmental predictors such as precipitation, vapor pressure deficit, land-cover composition, elevation, slope, and wind contain strong information for wildfire-loss modeling. Second, geographically stricter validation gives a more conservative and realistic estimate of model transferability than random splitting alone. Third, monetary wildfire-loss assessment can be complemented by ecological exposure screening, which helps identify counties where high predicted wildfire risk overlaps with high natural vegetation exposure. These insights make the framework transferable to other regions where comparable wildfire hazard, exposure, vulnerability, environmental, and loss-related datasets are available.

Overall, the thesis demonstrates that publicly accessible data can be used to construct a complete county-level wildfire risk prediction and assessment pipeline focused on monetary expected loss. The results support the broader thesis direction of moving from hazard-oriented wildfire prediction toward integrated, consequence-oriented wildfire-loss assessment.

% -------------------------
% Appendices
% -------------------------
\appendix
\include{appendices/appendix_supporting_outputs}
\include{appendices/appendix_creek_fire_pipeline}
\include{appendices/appendix_conference_certificate}

% -------------------------
% Bibliography
% -------------------------
\clearpage
\bibliography{references}

\end{document}